% Options for packages loaded elsewhere
\PassOptionsToPackage{unicode}{hyperref}
\PassOptionsToPackage{hyphens}{url}
\documentclass[
]{article}
\usepackage{xcolor}
\usepackage[margin=1in]{geometry}
\usepackage{amsmath,amssymb}
\setcounter{secnumdepth}{-\maxdimen} % remove section numbering
\usepackage{iftex}
\ifPDFTeX
  \usepackage[T1]{fontenc}
  \usepackage[utf8]{inputenc}
  \usepackage{textcomp} % provide euro and other symbols
\else % if luatex or xetex
  \usepackage{unicode-math} % this also loads fontspec
  \defaultfontfeatures{Scale=MatchLowercase}
  \defaultfontfeatures[\rmfamily]{Ligatures=TeX,Scale=1}
\fi
\usepackage{lmodern}
\ifPDFTeX\else
  % xetex/luatex font selection
\fi
% Use upquote if available, for straight quotes in verbatim environments
\IfFileExists{upquote.sty}{\usepackage{upquote}}{}
\IfFileExists{microtype.sty}{% use microtype if available
  \usepackage[]{microtype}
  \UseMicrotypeSet[protrusion]{basicmath} % disable protrusion for tt fonts
}{}
\makeatletter
\@ifundefined{KOMAClassName}{% if non-KOMA class
  \IfFileExists{parskip.sty}{%
    \usepackage{parskip}
  }{% else
    \setlength{\parindent}{0pt}
    \setlength{\parskip}{6pt plus 2pt minus 1pt}}
}{% if KOMA class
  \KOMAoptions{parskip=half}}
\makeatother
\usepackage{color}
\usepackage{fancyvrb}
\newcommand{\VerbBar}{|}
\newcommand{\VERB}{\Verb[commandchars=\\\{\}]}
\DefineVerbatimEnvironment{Highlighting}{Verbatim}{commandchars=\\\{\}}
% Add ',fontsize=\small' for more characters per line
\usepackage{framed}
\definecolor{shadecolor}{RGB}{248,248,248}
\newenvironment{Shaded}{\begin{snugshade}}{\end{snugshade}}
\newcommand{\AlertTok}[1]{\textcolor[rgb]{0.94,0.16,0.16}{#1}}
\newcommand{\AnnotationTok}[1]{\textcolor[rgb]{0.56,0.35,0.01}{\textbf{\textit{#1}}}}
\newcommand{\AttributeTok}[1]{\textcolor[rgb]{0.13,0.29,0.53}{#1}}
\newcommand{\BaseNTok}[1]{\textcolor[rgb]{0.00,0.00,0.81}{#1}}
\newcommand{\BuiltInTok}[1]{#1}
\newcommand{\CharTok}[1]{\textcolor[rgb]{0.31,0.60,0.02}{#1}}
\newcommand{\CommentTok}[1]{\textcolor[rgb]{0.56,0.35,0.01}{\textit{#1}}}
\newcommand{\CommentVarTok}[1]{\textcolor[rgb]{0.56,0.35,0.01}{\textbf{\textit{#1}}}}
\newcommand{\ConstantTok}[1]{\textcolor[rgb]{0.56,0.35,0.01}{#1}}
\newcommand{\ControlFlowTok}[1]{\textcolor[rgb]{0.13,0.29,0.53}{\textbf{#1}}}
\newcommand{\DataTypeTok}[1]{\textcolor[rgb]{0.13,0.29,0.53}{#1}}
\newcommand{\DecValTok}[1]{\textcolor[rgb]{0.00,0.00,0.81}{#1}}
\newcommand{\DocumentationTok}[1]{\textcolor[rgb]{0.56,0.35,0.01}{\textbf{\textit{#1}}}}
\newcommand{\ErrorTok}[1]{\textcolor[rgb]{0.64,0.00,0.00}{\textbf{#1}}}
\newcommand{\ExtensionTok}[1]{#1}
\newcommand{\FloatTok}[1]{\textcolor[rgb]{0.00,0.00,0.81}{#1}}
\newcommand{\FunctionTok}[1]{\textcolor[rgb]{0.13,0.29,0.53}{\textbf{#1}}}
\newcommand{\ImportTok}[1]{#1}
\newcommand{\InformationTok}[1]{\textcolor[rgb]{0.56,0.35,0.01}{\textbf{\textit{#1}}}}
\newcommand{\KeywordTok}[1]{\textcolor[rgb]{0.13,0.29,0.53}{\textbf{#1}}}
\newcommand{\NormalTok}[1]{#1}
\newcommand{\OperatorTok}[1]{\textcolor[rgb]{0.81,0.36,0.00}{\textbf{#1}}}
\newcommand{\OtherTok}[1]{\textcolor[rgb]{0.56,0.35,0.01}{#1}}
\newcommand{\PreprocessorTok}[1]{\textcolor[rgb]{0.56,0.35,0.01}{\textit{#1}}}
\newcommand{\RegionMarkerTok}[1]{#1}
\newcommand{\SpecialCharTok}[1]{\textcolor[rgb]{0.81,0.36,0.00}{\textbf{#1}}}
\newcommand{\SpecialStringTok}[1]{\textcolor[rgb]{0.31,0.60,0.02}{#1}}
\newcommand{\StringTok}[1]{\textcolor[rgb]{0.31,0.60,0.02}{#1}}
\newcommand{\VariableTok}[1]{\textcolor[rgb]{0.00,0.00,0.00}{#1}}
\newcommand{\VerbatimStringTok}[1]{\textcolor[rgb]{0.31,0.60,0.02}{#1}}
\newcommand{\WarningTok}[1]{\textcolor[rgb]{0.56,0.35,0.01}{\textbf{\textit{#1}}}}
\usepackage{longtable,booktabs,array}
\newcounter{none} % for unnumbered tables
\usepackage{calc} % for calculating minipage widths
% Correct order of tables after \paragraph or \subparagraph
\usepackage{etoolbox}
\makeatletter
\patchcmd\longtable{\par}{\if@noskipsec\mbox{}\fi\par}{}{}
\makeatother
% Allow footnotes in longtable head/foot
\IfFileExists{footnotehyper.sty}{\usepackage{footnotehyper}}{\usepackage{footnote}}
\makesavenoteenv{longtable}
\usepackage{graphicx}
\makeatletter
\newsavebox\pandoc@box
\newcommand*\pandocbounded[1]{% scales image to fit in text height/width
  \sbox\pandoc@box{#1}%
  \Gscale@div\@tempa{\textheight}{\dimexpr\ht\pandoc@box+\dp\pandoc@box\relax}%
  \Gscale@div\@tempb{\linewidth}{\wd\pandoc@box}%
  \ifdim\@tempb\p@<\@tempa\p@\let\@tempa\@tempb\fi% select the smaller of both
  \ifdim\@tempa\p@<\p@\scalebox{\@tempa}{\usebox\pandoc@box}%
  \else\usebox{\pandoc@box}%
  \fi%
}
% Set default figure placement to htbp
\def\fps@figure{htbp}
\makeatother
\setlength{\emergencystretch}{3em} % prevent overfull lines
\providecommand{\tightlist}{%
  \setlength{\itemsep}{0pt}\setlength{\parskip}{0pt}}
\usepackage{bookmark}
\IfFileExists{xurl.sty}{\usepackage{xurl}}{} % add URL line breaks if available
\urlstyle{same}
\hypersetup{
  pdftitle={review},
  hidelinks,
  pdfcreator={LaTeX via pandoc}}

\title{review}
\author{}
\date{\vspace{-2.5em}}

\begin{document}
\maketitle

\section{Review Packet: Forces, Equilibrium and Linear Motion (Chapters
2-3)}\label{review-packet-forces-equilibrium-and-linear-motion-chapters-2-3}

Oct 16, 2026

\subsection{How to use this packet}\label{how-to-use-this-packet}

This packet walks you through Chapters 2 and 3 in the order the ideas
build on each other. Every part has a short explanation, a worked
example, and practice questions, and the answer key is at the end.

Chapter 2 asks: \textbf{what makes motion change?} The answer is net
force. Chapter 3 asks: \textbf{how do we describe motion changing?} The
answer is velocity, acceleration, graphs, and a few equations. The two
chapters connect in one chain:

\[
\text{net force} = 0 \;\Longleftrightarrow\; \text{velocity does not change} \;\Longleftrightarrow\; \text{acceleration} = 0
\]

Keep that chain in mind. It comes back in almost every part.

\textbf{How to work through it}

\begin{itemize}
\tightlist
\item
  Read the \textbf{Big idea} first, then follow the \textbf{Worked
  example} with a pencil in your hand.
\item
  Try the practice questions before you look at the answers. A wrong
  attempt that you then fix teaches more than reading a right answer.
\item
  Show your work and write units: N, m, s, m/s, m/s².
\item
  Use g = 10 m/s² for falling objects unless a question says otherwise.
\item
  Questions marked ★ are challenges. Skip them on a first pass if you
  are short on time.
\end{itemize}

\textbf{Map of the packet}

{\def\LTcaptype{none} % do not increment counter
\begin{longtable}[]{@{}
  >{\raggedright\arraybackslash}p{(\linewidth - 4\tabcolsep) * \real{0.3333}}
  >{\raggedright\arraybackslash}p{(\linewidth - 4\tabcolsep) * \real{0.3333}}
  >{\raggedright\arraybackslash}p{(\linewidth - 4\tabcolsep) * \real{0.3333}}@{}}
\toprule\noalign{}
\begin{minipage}[b]{\linewidth}\raggedright
Part
\end{minipage} & \begin{minipage}[b]{\linewidth}\raggedright
Topic
\end{minipage} & \begin{minipage}[b]{\linewidth}\raggedright
The skill you will practice
\end{minipage} \\
\midrule\noalign{}
\endhead
\bottomrule\noalign{}
\endlastfoot
A1 & Inertia and Newton's first law & Explaining motion with no net
force \\
A2 & Force as a vector, net force & Adding forces along a line \\
A3 & Equilibrium & Static vs dynamic, finding an unknown force \\
A4 & Support force, tension, force diagrams & Reading and checking a
force diagram \\
A5 & Angles and tension, perpendicular forces & Rope angle reasoning,
Pythagorean theorem \\
B1 & Speed, velocity, relative motion & Average speed, adding speeds \\
B2 & Position-time graphs & Slope = velocity \\
B3 & Tables, acceleration & Estimating velocity and acceleration from
data \\
B4 & Velocity-time graphs & Slope = acceleration \\
B5 & Constant acceleration equations & Choosing and using v = at and d =
½at² \\
B6 & Free fall and tosses & g, symmetry, the top of the path \\
B7 & Adding velocity vectors & Tip-to-tail diagrams, crosswinds and
currents \\
C & Putting it together & Linking forces and motion, then a full mixed
test \\
\end{longtable}
}

\subsection{Reference sheet: formulas and when to use
them}\label{reference-sheet-formulas-and-when-to-use-them}

Every formula you need for Chapters 2 and 3 is in the tables below,
along with the situation that tells you to use it. Look here whenever
you are not sure where to start.

\textbf{Forces (Chapter 2)}

{\def\LTcaptype{none} % do not increment counter
\begin{longtable}[]{@{}
  >{\raggedright\arraybackslash}p{(\linewidth - 4\tabcolsep) * \real{0.3333}}
  >{\raggedright\arraybackslash}p{(\linewidth - 4\tabcolsep) * \real{0.3333}}
  >{\raggedright\arraybackslash}p{(\linewidth - 4\tabcolsep) * \real{0.3333}}@{}}
\toprule\noalign{}
\begin{minipage}[b]{\linewidth}\raggedright
Tool
\end{minipage} & \begin{minipage}[b]{\linewidth}\raggedright
What it says
\end{minipage} & \begin{minipage}[b]{\linewidth}\raggedright
Use it when
\end{minipage} \\
\midrule\noalign{}
\endhead
\bottomrule\noalign{}
\endlastfoot
Net force on a line & Add forces in the same direction; subtract forces
in opposite directions. Call right/up positive, left/down negative. &
Several forces push or pull along the same line and you want the
total. \\
Equilibrium: ΣF = 0 & Net force is zero. The object is at rest (static)
or moving at constant velocity (dynamic). & You are told an object is at
rest or at constant velocity, or you need an unknown force. \\
Weight: W = m × g & Weight in newtons = mass in kg × 10 N/kg (on Earth).
& You know the mass and need the downward pull of gravity. \\
Support force & For an object at rest on a level surface with no other
vertical forces, support force = weight. & A book on a table, a person
on a scale. \\
Tension & A stretching force along a rope. A rope pulled equally from
both ends has a tension equal to ONE end's pull. & Tug-of-war, ropes
tied to trees, hanging objects. \\
Rope angle & The flatter the rope (smaller angle above horizontal), the
larger the tension needed to hold the same weight. & Hammocks,
clotheslines, hanging signs. \\
Right-angle forces: R = √(a² + b²) & The resultant of two perpendicular
forces (the Pythagorean theorem). & Two forces at 90° to each other. \\
Tip-to-tail rule & Draw each force starting where the last one ended. If
the arrows return to the start, the net force is zero. & Showing that
forces balance on a diagram. \\
\end{longtable}
}

\textbf{Motion (Chapter 3)}

{\def\LTcaptype{none} % do not increment counter
\begin{longtable}[]{@{}
  >{\raggedright\arraybackslash}p{(\linewidth - 4\tabcolsep) * \real{0.3333}}
  >{\raggedright\arraybackslash}p{(\linewidth - 4\tabcolsep) * \real{0.3333}}
  >{\raggedright\arraybackslash}p{(\linewidth - 4\tabcolsep) * \real{0.3333}}@{}}
\toprule\noalign{}
\begin{minipage}[b]{\linewidth}\raggedright
Tool
\end{minipage} & \begin{minipage}[b]{\linewidth}\raggedright
What it says
\end{minipage} & \begin{minipage}[b]{\linewidth}\raggedright
Use it when
\end{minipage} \\
\midrule\noalign{}
\endhead
\bottomrule\noalign{}
\endlastfoot
Average speed = distance ÷ time & How fast, with no direction. & You
know how far and how long. \\
d = v × t & Distance at a steady speed. & Speed is constant, or you are
using an average speed or average velocity. \\
v ≈ Δx ÷ Δt & Velocity is the change in position divided by the change
in time. & You have two position readings close together in time, or the
slope of a position-time graph. \\
a = Δv ÷ Δt & Acceleration = (final velocity − starting velocity) ÷ time
taken. Units m/s². & You know two velocities and the time between them,
or the slope of a velocity-time graph. \\
v = v₀ + a × t & Velocity after time t. If the object starts from rest,
v₀ = 0 and this is just v = a × t. & Constant acceleration and you want
the speed at a certain time. \\
v\_avg = (v₀ + v) ÷ 2 & Average velocity when acceleration is constant.
& Constant acceleration and you want distance without the next
formula. \\
d = ½ × a × t² & Distance covered from rest at constant acceleration. &
Starting from rest and you know a and t. \\
Free fall & a = g = 10 m/s² downward. Dropped from rest: v = g × t and d
= ½ × g × t² = 5 × t². & Only gravity acts (ignore air resistance),
going up, down, or at the top. \\
Toss straight up & Time up = v₀ ÷ g. Total time in air = 2 × time up.
Return speed = launch speed. Max height = distance fallen in the time
up. & A ball thrown up and coming back to the same level. \\
Adding velocities & Ground velocity = object's velocity + medium's
velocity. Same direction: add. Opposite: subtract. Perpendicular: R =
√(a² + b²). & Boats in currents, planes in wind, walking on moving
trains. \\
\end{longtable}
}

\textbf{Reading graphs}

{\def\LTcaptype{none} % do not increment counter
\begin{longtable}[]{@{}
  >{\raggedright\arraybackslash}p{(\linewidth - 4\tabcolsep) * \real{0.3333}}
  >{\raggedright\arraybackslash}p{(\linewidth - 4\tabcolsep) * \real{0.3333}}
  >{\raggedright\arraybackslash}p{(\linewidth - 4\tabcolsep) * \real{0.3333}}@{}}
\toprule\noalign{}
\begin{minipage}[b]{\linewidth}\raggedright
Feature
\end{minipage} & \begin{minipage}[b]{\linewidth}\raggedright
On a position-time graph
\end{minipage} & \begin{minipage}[b]{\linewidth}\raggedright
On a velocity-time graph
\end{minipage} \\
\midrule\noalign{}
\endhead
\bottomrule\noalign{}
\endlastfoot
Height of the line & Where the object is & How fast and which way it
moves \\
Slope of the line & \textbf{Velocity} & \textbf{Acceleration} \\
Flat (horizontal) line & At rest & Constant velocity, acceleration =
0 \\
Straight sloped line & Constant velocity & Constant acceleration \\
Slope up / slope down & Moving in the positive / negative direction &
Speeding up / slowing down in the positive direction \\
Line crosses zero & Object passes the origin & Object turns around
(velocity changes sign) \\
\end{longtable}
}

\textbf{Falling objects from rest (g = 10 m/s²)}

{\def\LTcaptype{none} % do not increment counter
\begin{longtable}[]{@{}lll@{}}
\toprule\noalign{}
Time (s) & Speed (m/s) & Total distance fallen (m) \\
\midrule\noalign{}
\endhead
\bottomrule\noalign{}
\endlastfoot
0 & 0 & 0 \\
1 & 10 & 5 \\
2 & 20 & 20 \\
3 & 30 & 45 \\
4 & 40 & 80 \\
5 & 50 & 125 \\
\end{longtable}
}

\subsection{Part A1: Inertia and Newton's first
law}\label{part-a1-inertia-and-newtons-first-law}

\textbf{Big idea.} Objects keep doing what they are already doing unless
a net force changes it. An object at rest stays at rest, and an object
in motion keeps moving at the same speed in a straight line. This
tendency is called \textbf{inertia}. Inertia is a property of matter,
not a force, and an object with more mass has more inertia.

This was a big change in thinking. Aristotle taught that moving objects
need a continual push to keep moving. Galileo, using ramps and rolling
balls, realized that things slow down because of \textbf{friction}, a
force, and not because stopping is their ``natural'' state. Remove
friction and a moving object just keeps going.

\textbf{Watch out:} ``No net force'' does not mean ``no motion.'' It
means \emph{no change} in motion.

\textbf{Inertia in everyday life}

\begin{itemize}
\tightlist
\item
  Yank a tablecloth out from under dishes and the dishes stay put. They
  were at rest, and the quick tug gave them almost no push.
\item
  You are standing on a bus that turns left, and you lurch to the right.
  Your body was moving straight ahead and keeps trying to. The bus
  turned out from under you.
\item
  Toss a coin straight up on a train moving at constant velocity and it
  lands in your hand. The coin shares the train's forward motion and
  nothing pushes it forward or backward.
\item
  Toss the coin while the train speeds up and it lands \emph{behind}
  you. The coin keeps its old forward velocity while the train moves
  ahead.
\end{itemize}

\textbf{Worked example.} A seagull flies in a straight line at a
constant 10 m/s and lets go of a clam over the rocks, hoping the fall
will crack the shell. Ignoring air resistance, how fast is the clam
moving horizontally just before it hits the rocks?

\begin{enumerate}
\def\labelenumi{\arabic{enumi}.}
\tightlist
\item
  Before the gull let go, the clam moved forward at 10 m/s along with
  the gull.
\item
  Once released, gravity pulls the clam down, but nothing pushes it
  forward or backward.
\item
  No horizontal force means no change in horizontal velocity. The answer
  is \textbf{10 m/s}. Seen from the ground, the clam keeps moving
  forward as it falls, so it lands well ahead of the spot directly below
  where it was released. Seen from the gull, which keeps flying at the
  same velocity, the clam stays directly below the bird the whole way
  down.
\end{enumerate}

\textbf{Try it}

\textbf{A1.1} If no external forces act on a moving object, it will

\begin{itemize}
\tightlist
\item
  \begin{enumerate}
  \def\labelenumi{\Alph{enumi})}
  \tightlist
  \item
    continue moving at the same speed in a straight line.
  \end{enumerate}
\item
  \begin{enumerate}
  \def\labelenumi{\Alph{enumi})}
  \setcounter{enumi}{1}
  \tightlist
  \item
    slow down gradually until it stops.
  \end{enumerate}
\item
  \begin{enumerate}
  \def\labelenumi{\Alph{enumi})}
  \setcounter{enumi}{2}
  \tightlist
  \item
    stop abruptly.
  \end{enumerate}
\item
  \begin{enumerate}
  \def\labelenumi{\Alph{enumi})}
  \setcounter{enumi}{3}
  \tightlist
  \item
    speed up.
  \end{enumerate}
\end{itemize}

\textbf{A1.2} A rock drifts through deep space at 20 m/s, far from any
planet. What force is needed to keep it moving at 20 m/s?

\begin{itemize}
\tightlist
\item
  \begin{enumerate}
  \def\labelenumi{\Alph{enumi})}
  \tightlist
  \item
    A force equal to its weight
  \end{enumerate}
\item
  \begin{enumerate}
  \def\labelenumi{\Alph{enumi})}
  \setcounter{enumi}{1}
  \tightlist
  \item
    20 N
  \end{enumerate}
\item
  \begin{enumerate}
  \def\labelenumi{\Alph{enumi})}
  \setcounter{enumi}{2}
  \tightlist
  \item
    No force at all
  \end{enumerate}
\item
  \begin{enumerate}
  \def\labelenumi{\Alph{enumi})}
  \setcounter{enumi}{3}
  \tightlist
  \item
    It depends on the rock's mass
  \end{enumerate}
\end{itemize}

\textbf{A1.3} You whirl a rock on a string in a horizontal circle, and
the string suddenly breaks. Seen from above, the rock then moves

\begin{itemize}
\tightlist
\item
  \begin{enumerate}
  \def\labelenumi{\Alph{enumi})}
  \tightlist
  \item
    in a circle, as before.
  \end{enumerate}
\item
  \begin{enumerate}
  \def\labelenumi{\Alph{enumi})}
  \setcounter{enumi}{1}
  \tightlist
  \item
    in a straight line in the direction it was moving at the moment the
    string broke.
  \end{enumerate}
\item
  \begin{enumerate}
  \def\labelenumi{\Alph{enumi})}
  \setcounter{enumi}{2}
  \tightlist
  \item
    in a straight line directly away from the center of the circle.
  \end{enumerate}
\item
  \begin{enumerate}
  \def\labelenumi{\Alph{enumi})}
  \setcounter{enumi}{3}
  \tightlist
  \item
    toward the center of the circle.
  \end{enumerate}
\end{itemize}

\textbf{A1.4} Which has more inertia, a bowling ball or a tennis ball?
How do you know?

\textbf{A1.5} A hockey puck slides across ice and slowly comes to a
stop. Explain \emph{why} it stops, and explain what would happen on
perfectly frictionless ice.

\textbf{A1.6} A coin is tossed straight up inside a train moving at
constant velocity. Explain why the coin lands back in your hand. Then
explain where it lands if the train speeds up while the coin is in the
air.

\subsection{Part A2: Force as a vector and finding net
force}\label{part-a2-force-as-a-vector-and-finding-net-force}

\textbf{Big idea.} A force is a push or pull, measured in
\textbf{newtons (N)}. A force is a \textbf{vector}, meaning it has both
a \textbf{magnitude} (how strong) and a \textbf{direction} (which way).
You need both to describe what a force does.

When several forces act on an object, what matters is the \textbf{net
force}, the single combined force. Whether the velocity changes depends
on the net force:

\begin{itemize}
\tightlist
\item
  \textbf{Net force = 0:} the velocity does not change.
\item
  \textbf{Net force ≠ 0:} the velocity changes. The speed can change,
  the direction can change, or both.
\end{itemize}

\textbf{How to find net force along a line}

\begin{enumerate}
\def\labelenumi{\arabic{enumi}.}
\tightlist
\item
  Pick a positive direction (for example right is +, left is −; up is +,
  down is −).
\item
  Add the forces, using the signs.
\item
  The sign of the answer tells you the direction. The size is the
  magnitude.
\end{enumerate}

Forces at right angles to each other do \textbf{not} cancel. Handle each
direction on its own: left-right forces with left-right forces, up-down
forces with up-down forces.

\textbf{Worked example 1.} A block is pulled 5 N to the left and 15 N to
the right. Right is positive:

\begin{Shaded}
\begin{Highlighting}[]
\FunctionTok{\textbackslash{}text}\NormalTok{\{net\} = (+15) + ({-}5) = +10 }\FunctionTok{\textbackslash{}text}\NormalTok{\{ N\}}
\end{Highlighting}
\end{Shaded}

The net force is \textbf{10 N to the right}.

\textbf{Worked example 2.} An object feels 120 N to the left, 70 N to
the right, 45 N up, and 45 N down.

\begin{itemize}
\tightlist
\item
  Up-down: +45 − 45 = 0 N, so these cancel.
\item
  Left-right: +70 − 120 = −50 N.
\end{itemize}

The net force is \textbf{50 N to the left}.

\textbf{Try it}

\textbf{A2.1} A tree stump is pulled north by a 10 N force and south by
a 25 N force. What is the net force (magnitude and direction)?

\textbf{A2.2} A box is pushed 30 N to the right, another 20 N to the
right, and 10 N to the left. What is the net force?

\textbf{A2.3} An object has these forces on it: 100 N up, 60 N down, 35
N left, and 35 N right. What is the net force?

\textbf{A2.4} True or false, with a reason: ``If the net force on an
object is zero, the object must be at rest.''

\textbf{A2.5} Why is force called a \emph{vector} quantity?

\textbf{A2.6} Two 10 N forces act on a box. Describe the directions of
the two forces that would make the net force (a) 20 N and (b) 0 N.

\textbf{A2.7} ★ A cart is moving to the right. A net force of 8 N points
to the \emph{left}. Without any calculation, describe what happens to
the cart's velocity over the next few seconds.

\subsection{Part A3: Equilibrium}\label{part-a3-equilibrium}

\textbf{Big idea.} An object is in \textbf{equilibrium} when the net
force on it is zero. There are two kinds, and both have balanced forces:

\begin{itemize}
\tightlist
\item
  \textbf{Static equilibrium:} the object is at rest.
\item
  \textbf{Dynamic equilibrium:} the object moves at constant velocity
  (constant speed in a straight line).
\end{itemize}

\begin{Shaded}
\begin{Highlighting}[]
\FunctionTok{\textbackslash{}Sigma}\NormalTok{ F = 0}
\end{Highlighting}
\end{Shaded}

\textbf{Watch out:} Many people think ``it is moving, so there must be a
net force.'' That is Aristotle's idea, and Galileo showed it is wrong.
Constant velocity means zero net force. Only a \emph{change} in velocity
needs a net force.

A falling insect that reaches a steady speed is a good example. Gravity
pulls it down with its weight, and air resistance pushes up with an
equal force. The forces balance, so the speed stops changing.

\textbf{Finding an unknown force.} If an object is in equilibrium, the
known forces and the unknown force must add to zero. So the unknown
force is equal in size and opposite in direction to the sum of the known
forces.

\textbf{Worked example.} Three horizontal forces act on an object in
static equilibrium. Two of them are 80 N to the right and 30 N to the
left, and the third is unknown. Find the unknown force.

\begin{enumerate}
\def\labelenumi{\arabic{enumi}.}
\tightlist
\item
  Right is positive. The known forces add to +80 − 30 = +50 N, which is
  50 N to the right.
\item
  For zero net force, the unknown force must cancel that: it is
  \textbf{50 N to the left}.
\item
  Check: +80 − 30 − 50 = 0.
\end{enumerate}

\textbf{Try it}

\textbf{A3.1} A book sits motionless on a table. Which best describes
it?

\begin{itemize}
\tightlist
\item
  \begin{enumerate}
  \def\labelenumi{\Alph{enumi})}
  \tightlist
  \item
    Static equilibrium
  \end{enumerate}
\item
  \begin{enumerate}
  \def\labelenumi{\Alph{enumi})}
  \setcounter{enumi}{1}
  \tightlist
  \item
    Dynamic equilibrium
  \end{enumerate}
\item
  \begin{enumerate}
  \def\labelenumi{\Alph{enumi})}
  \setcounter{enumi}{2}
  \tightlist
  \item
    Not in equilibrium, because gravity pulls it down
  \end{enumerate}
\item
  \begin{enumerate}
  \def\labelenumi{\Alph{enumi})}
  \setcounter{enumi}{3}
  \tightlist
  \item
    No forces act on it
  \end{enumerate}
\end{itemize}

\textbf{A3.2} A hockey puck glides at constant velocity across smooth
ice. Is it in equilibrium? If so, which kind? What is the net force on
it?

\textbf{A3.3} An object is in static equilibrium. Three horizontal
forces act on it: 120 N to the right, 45 N to the left, and an unknown
force. Find the magnitude and direction of the unknown force.

\textbf{A3.4} A parachutist who weighs 600 N falls at a constant
velocity. What is the net force on her? What is the air resistance,
including its direction?

\textbf{A3.5} A car drives at a steady 25 m/s on a straight, level road.
The engine pushes the car forward with 2000 N. What is the total
friction and air resistance on the car, and what kind of equilibrium is
this?

\textbf{A3.6} A friend says, ``The box is sliding across the floor, so
there must be a net force on it.'' Explain whether your friend is right.

\textbf{A3.7} ★ An object has these forces on it: 50 N up, 50 N down, 30
N left, and 30 N right. A classmate says it must be at rest. What is
wrong with that claim?

\subsection{Part A4: Support force, tension, and force
diagrams}\label{part-a4-support-force-tension-and-force-diagrams}

\textbf{Big idea.} Four forces show up again and again, and a
\textbf{force diagram} is how you keep track of them. The diagram shows
the object as a dot, with an arrow for each force acting on it. Arrow
length is proportional to the force's size, and each arrow is labeled.

{\def\LTcaptype{none} % do not increment counter
\begin{longtable}[]{@{}
  >{\raggedright\arraybackslash}p{(\linewidth - 4\tabcolsep) * \real{0.3333}}
  >{\raggedright\arraybackslash}p{(\linewidth - 4\tabcolsep) * \real{0.3333}}
  >{\raggedright\arraybackslash}p{(\linewidth - 4\tabcolsep) * \real{0.3333}}@{}}
\toprule\noalign{}
\begin{minipage}[b]{\linewidth}\raggedright
Force
\end{minipage} & \begin{minipage}[b]{\linewidth}\raggedright
What it is
\end{minipage} & \begin{minipage}[b]{\linewidth}\raggedright
Direction
\end{minipage} \\
\midrule\noalign{}
\endhead
\bottomrule\noalign{}
\endlastfoot
\textbf{Weight} & Gravity's pull on the object. W = m × g, with m in kg
and g = 10 N/kg. & Always down \\
\textbf{Support (normal) force} & The push of a surface on the object. &
Perpendicular to the surface, away from it \\
\textbf{Tension} & A stretching pull along a rope or string. & Along the
rope, away from the object \\
\textbf{Friction} & A force from a surface that resists sliding. &
Opposite to the motion (or the attempted motion) \\
\end{longtable}
}

\textbf{Key facts}

\begin{itemize}
\tightlist
\item
  A book at rest on a table has a support force equal to its weight, so
  the two cancel.
\item
  A scale reads the support force. If you stand on two scales, the two
  readings \textbf{add up} to your weight.
\item
  Tension is a force that stretches. If two people pull on opposite ends
  of a rope with 400 N each and nothing moves, the net force on the rope
  is \textbf{0 N} and the tension is \textbf{400 N}, not 800 N. Tension
  equals one end's pull.
\item
  If a scaffold hangs from ropes at rest, the tensions in the ropes
  \textbf{add up} to the total weight of the scaffold plus the people on
  it.
\end{itemize}

\textbf{Tip-to-tail check.} If an object is in equilibrium, the force
arrows drawn tip-to-tail form a closed shape that comes back to the
starting point. If the shape does not close, there is a net force.

\textbf{Worked example.} A sled moves at constant velocity as a person
pulls it with a rope that points up and to the right. The forces are:

\begin{itemize}
\tightlist
\item
  Weight: 400 N, down.
\item
  Support force: 350 N, up.
\item
  Tension: 130 N along the rope. Its horizontal part is 120 N to the
  right and its vertical part is 50 N up.
\item
  Friction: 120 N, to the left.
\end{itemize}

Show that the sled is in dynamic equilibrium.

\begin{enumerate}
\def\labelenumi{\arabic{enumi}.}
\tightlist
\item
  \textbf{Vertical forces:} 350 N up + 50 N up (from the rope) = 400 N
  up, which balances the 400 N weight down. Net vertical force = 0.
\item
  \textbf{Horizontal forces:} 120 N right (from the rope) balances 120 N
  left (friction). Net horizontal force = 0.
\item
  Tip-to-tail: start at a point, go down 400, up 350, then along the
  rope (right 120, up 50), then left 120. You are back at the start, so
  the forces balance.
\item
  Net force = 0 and the sled moves at constant velocity, so it is in
  \textbf{dynamic equilibrium}.
\end{enumerate}

\textbf{Try it}

\textbf{A4.1} A 10 N book rests on a table. What is the support force on
it?

\textbf{A4.2} A 30 kg dog is sleeping on the floor. What is its weight,
and what is the support force on it?

\textbf{A4.3} A 500 N student stands at rest on two bathroom scales,
with her weight shared equally. What does each scale read?

\textbf{A4.4} Two painters with a combined weight of 1300 N stand on a
scaffold that weighs 400 N. The scaffold hangs at rest from two ropes.
What do the two rope tensions add up to?

\textbf{A4.5} In a tug-of-war that is perfectly balanced, each team
pulls on the rope with 300 N. What is the net force on the rope? What is
the tension in the rope?

\textbf{A4.6} A crate is pulled across the floor at constant velocity by
a rope. The forces are: weight 800 N down, support force 650 N up, and
tension whose horizontal part is 240 N to the right and whose vertical
part is 150 N up. (a) What is the friction force? (b) Show that the
vertical forces balance.

\textbf{A4.7} A person weighs themself on a bathroom scale, then rises
up onto their tiptoes and stays still. Does the scale reading change?
Explain using equilibrium.

\textbf{A4.8} ★ In the worked example above, the person pulls harder and
the sled's speed starts increasing. What does that tell you about the
horizontal forces?

\subsection{Part A5: Angles and tension, and forces at right
angles}\label{part-a5-angles-and-tension-and-forces-at-right-angles}

\textbf{Big idea, part 1: rope angles.} When an object hangs at rest
from two ropes, the \textbf{upward parts} of the two tensions have to
add up to the object's weight. Each rope only contributes the part of
its pull that points up.

\begin{itemize}
\tightlist
\item
  If both ropes are \textbf{vertical}, all of each tension points up.
  Each rope carries \textbf{half the weight}.
\item
  If the ropes are \textbf{tilted} away from vertical, only part of each
  tension points up, so each rope must pull harder. Each tension is
  \textbf{more than half the weight}.
\item
  The \textbf{flatter} the ropes (the smaller the angle above
  horizontal), the less each rope's pull points up, and the
  \textbf{larger the tension} has to be.
\end{itemize}

Tension depends on the \textbf{angle} of the ropes, not on their length.
If both ropes make equal angles, the two tensions are equal. If the
angles are different, the tensions are different.

\textbf{Hammock example.} A hammock hangs between two trees. When the
trees are far apart, the ropes are nearly flat, the angle of elevation
is small, and the tension is large. When the trees are closer together,
the ropes tilt more steeply, the angle is larger, and the tension is
smaller.

{\def\LTcaptype{none} % do not increment counter
\begin{longtable}[]{@{}
  >{\raggedright\arraybackslash}p{(\linewidth - 4\tabcolsep) * \real{0.3333}}
  >{\raggedright\arraybackslash}p{(\linewidth - 4\tabcolsep) * \real{0.3333}}
  >{\raggedright\arraybackslash}p{(\linewidth - 4\tabcolsep) * \real{0.3333}}@{}}
\toprule\noalign{}
\begin{minipage}[b]{\linewidth}\raggedright
Ropes
\end{minipage} & \begin{minipage}[b]{\linewidth}\raggedright
Angle above horizontal
\end{minipage} & \begin{minipage}[b]{\linewidth}\raggedright
Tension in each rope
\end{minipage} \\
\midrule\noalign{}
\endhead
\bottomrule\noalign{}
\endlastfoot
Straight up (vertical) & 90° & Smallest: half the weight \\
Tilted & Medium & Greater than half the weight \\
Nearly flat & Small & Very large \\
Perfectly flat & 0° & Impossible, because no part of the pull points
up \\
\end{longtable}
}

\textbf{Big idea, part 2: forces at right angles.} Two forces that point
in perpendicular directions do not simply add. Their resultant is the
hypotenuse of a right triangle:

\begin{Shaded}
\begin{Highlighting}[]
\NormalTok{R = }\FunctionTok{\textbackslash{}sqrt}\NormalTok{\{a\^{}2 + b\^{}2\}}
\end{Highlighting}
\end{Shaded}

The same formula works for velocities at right angles (Part B7). Useful
sets of numbers: 3-4-5, 5-12-13, 6-8-10, and two equal forces at right
angles give about 1.4 times one force (10 N right and 10 N up give about
14 N right-up).

\textbf{Worked example.} A 40 N force pulls an object east and a 30 N
force pulls it north at the same time. What is the net force?

\begin{enumerate}
\def\labelenumi{\arabic{enumi}.}
\tightlist
\item
  The forces are at right angles, so use the Pythagorean theorem.
\item
  R = √(40² + 30²) = √(1600 + 900) = √2500 = \textbf{50 N}.
\item
  The direction is between east and north, tilted a bit more toward
  east, since the east force is larger.
\end{enumerate}

\textbf{Try it}

\textbf{A5.1} A monkey hangs from the middle of a horizontal rope. The
tension in the rope is greatest when the rope

\begin{itemize}
\tightlist
\item
  \begin{enumerate}
  \def\labelenumi{\Alph{enumi})}
  \tightlist
  \item
    sags a lot.
  \end{enumerate}
\item
  \begin{enumerate}
  \def\labelenumi{\Alph{enumi})}
  \setcounter{enumi}{1}
  \tightlist
  \item
    sags very little.
  \end{enumerate}
\item
  \begin{enumerate}
  \def\labelenumi{\Alph{enumi})}
  \setcounter{enumi}{2}
  \tightlist
  \item
    has the same tension no matter how much it sags.
  \end{enumerate}
\end{itemize}

\textbf{A5.2} An 800 N object hangs at rest from two ropes that are both
vertical. What is the tension in each rope?

\textbf{A5.3} The same 800 N object is now hung from two ropes that are
tilted away from vertical, at equal angles. Is the tension in each rope
less than, equal to, or greater than 400 N? Explain.

\textbf{A5.4} You move the two trees holding a hammock farther apart,
but the hammock sags to the same spot. What happens to the angle of the
ropes and to the tension in them?

\textbf{A5.5} A 6 N force pulls an object east and an 8 N force pulls it
north. What is the magnitude of the net force?

\textbf{A5.6} Two 10 N forces act on a box at right angles. Is the net
force closer to 10 N, 14 N, or 20 N?

\textbf{A5.7} A 5 N force pulls a cart east while a 12 N force pulls it
north. Find the net force magnitude.

\textbf{A5.8} ★ A person tries to pull a clothesline completely flat and
horizontal while a heavy shirt hangs from its middle. Use what you know
about the upward parts of the tensions to explain why this can never be
done.

\subsection{Part B1: Speed, velocity, and relative
motion}\label{part-b1-speed-velocity-and-relative-motion}

\textbf{Big idea.} \textbf{Speed} tells you how fast something moves.
\textbf{Velocity} tells you how fast \emph{and} in which direction, so
velocity is a vector and speed is not. Both are measured in units such
as m/s or km/h.

\begin{Shaded}
\begin{Highlighting}[]
\FunctionTok{\textbackslash{}text}\NormalTok{\{average speed\} = }\FunctionTok{\textbackslash{}frac}\NormalTok{\{}\FunctionTok{\textbackslash{}text}\NormalTok{\{distance\}\}\{}\FunctionTok{\textbackslash{}text}\NormalTok{\{time\}\} }\FunctionTok{\textbackslash{}qquad\textbackslash{}qquad}\NormalTok{ d = v }\FunctionTok{\textbackslash{}times}\NormalTok{ t}
\end{Highlighting}
\end{Shaded}

\begin{itemize}
\tightlist
\item
  \textbf{Average speed} covers a whole trip: total distance divided by
  total time.
\item
  \textbf{Instantaneous speed} is the speed at one moment. A car's
  speedometer reads instantaneous speed.
\item
  A velocity of +5 m/s and a velocity of −5 m/s have the same speed but
  opposite directions.
\end{itemize}

\textbf{Relative motion.} Speeds are always measured relative to
something. When something moves on top of something else that is also
moving, the velocities combine:

\begin{itemize}
\tightlist
\item
  Moving in the \textbf{same} direction: \textbf{add} the speeds.
\item
  Moving in \textbf{opposite} directions: \textbf{subtract} the speeds.
\end{itemize}

\textbf{Worked example 1.} A horse gallops 10 km in 30 minutes. Convert
30 minutes to 0.5 hours:

\begin{Shaded}
\begin{Highlighting}[]
\FunctionTok{\textbackslash{}text}\NormalTok{\{average speed\} = }\FunctionTok{\textbackslash{}frac}\NormalTok{\{10 }\FunctionTok{\textbackslash{}text}\NormalTok{\{ km\}\}\{0.5 }\FunctionTok{\textbackslash{}text}\NormalTok{\{ h\}\} = 20 }\FunctionTok{\textbackslash{}text}\NormalTok{\{ km/h\}}
\end{Highlighting}
\end{Shaded}

\textbf{Worked example 2.} Jake runs at 4 m/s along a train flatcar that
moves at 10 m/s. Relative to the ground:

\begin{itemize}
\tightlist
\item
  Running in the same direction as the train: 10 + 4 = \textbf{14 m/s}.
\item
  Running in the opposite direction: 10 − 4 = \textbf{6 m/s}.
\end{itemize}

\textbf{Try it}

\textbf{B1.1} A car's speedometer shows how fast the car is going at one
moment. Is that average speed or instantaneous speed?

\textbf{B1.2} A runner covers 400 m in 80 s. What is her average speed?

\textbf{B1.3} A cyclist averages 6 m/s for 2 minutes. How far does the
cyclist travel? (Convert the time to seconds first.)

\textbf{B1.4} A car travels 60 km in 45 minutes. What is its average
speed in km/h?

\textbf{B1.5} Which of these are vectors: speed, velocity, distance,
acceleration?

\textbf{B1.6} Katelyn runs along the aisle of a train moving at 8 m/s.
Her speed relative to the floor of the train is 3 m/s. What could her
speed be relative to the ground? Give both possibilities and say when
each happens.

\textbf{B1.7} A mosquito flies at 3 m/s straight into a 3 m/s breeze
that blows in the opposite direction. What is the mosquito's speed
relative to the ground?

\textbf{B1.8} A car drives around a circular track at a constant 20 m/s.
Is its \emph{speed} constant? Is its \emph{velocity} constant? Explain.

\subsection{Part B2: Position-time
graphs}\label{part-b2-position-time-graphs}

\textbf{Big idea.} A position-time graph shows \textbf{where} an object
is at each moment. The \textbf{slope} of the graph is the object's
\textbf{velocity}:

\begin{Shaded}
\begin{Highlighting}[]
\NormalTok{v = }\FunctionTok{\textbackslash{}text}\NormalTok{\{slope\} = }\FunctionTok{\textbackslash{}frac}\NormalTok{\{}\FunctionTok{\textbackslash{}Delta}\NormalTok{ x\}\{}\FunctionTok{\textbackslash{}Delta}\NormalTok{ t\} = }\FunctionTok{\textbackslash{}frac}\NormalTok{\{}\FunctionTok{\textbackslash{}text}\NormalTok{\{rise\}\}\{}\FunctionTok{\textbackslash{}text}\NormalTok{\{run\}\}}
\end{Highlighting}
\end{Shaded}

\begin{itemize}
\tightlist
\item
  A \textbf{steeper} line means a \textbf{faster} object.
\item
  A line sloping \textbf{up} means motion in the \textbf{positive}
  direction. A line sloping \textbf{down} means motion in the
  \textbf{negative} direction.
\item
  A \textbf{flat} (horizontal) line means the object is \textbf{at
  rest}, because its position is not changing.
\end{itemize}

\textbf{Watch out:} A position-time graph is not a picture of the path.
A flat line does not mean ``moving at a constant speed''. It means
\emph{not moving}. Also, a line that is high on the graph just means the
object is far from the origin, not that it moves fast.

\textbf{Figure 1} shows a student walking back and forth along a
straight hallway. Position is measured from the doorway (x = 0), and
positive positions are toward the east end of the hall.

{[}embedded content: Practice data written for this packet; each dot is
a corner of the graph at whole-number values{]}

\textbf{Worked example.} Find the student's velocity from t = 0 to t = 4
s in Figure 1.

\begin{enumerate}
\def\labelenumi{\arabic{enumi}.}
\tightlist
\item
  At t = 0 the position is 0 m. At t = 4 s the position is 8 m.
\item
  Slope = rise ÷ run = (8 − 0) ÷ (4 − 0) = 8 ÷ 4 = \textbf{+2 m/s}.
\item
  The sign is positive, so the student is walking toward the east end at
  2 m/s.
\end{enumerate}

The straight line means the velocity is constant during that interval.
To find the velocity at t = 2 s, you do not need a new calculation: it
is the same, +2 m/s.

\textbf{Try it} (use Figure 1)

\textbf{B2.1} What is the student's position at t = 6 s?

\textbf{B2.2} What is the student's position at t = 10 s?

\textbf{B2.3} What is the student's velocity (speed and direction) at t
= 6 s?

\textbf{B2.4} What is the student's velocity at t = 10 s?

\textbf{B2.5} What is the student's velocity at t = 18 s?

\textbf{B2.6} During which time interval(s) is the student at rest?

\textbf{B2.7} During which time interval(s) is the student east of the
doorway (position greater than 0)?

\textbf{B2.8} During which time interval(s) is the student moving in the
positive direction?

\textbf{B2.9} During which interval is the student moving fastest, and
how fast?

\textbf{B2.10} ★ At t = 10 s the student is exactly at the doorway (x =
0). Is the student at rest at that moment? Explain using the slope.

\textbf{B2.11} ★ What is the total \emph{distance} the student walks
from t = 0 to t = 20 s? How does it compare with the student's final
\emph{position}?

\subsection{Part B3: Velocity from data tables, and
acceleration}\label{part-b3-velocity-from-data-tables-and-acceleration}

\textbf{Big idea.} \textbf{Acceleration} is how quickly velocity
changes. Any change in velocity counts: speeding up, slowing down, or
changing direction. Acceleration is a vector and its units are m/s²
(meters per second, every second).

\begin{Shaded}
\begin{Highlighting}[]
\NormalTok{a = }\FunctionTok{\textbackslash{}frac}\NormalTok{\{}\FunctionTok{\textbackslash{}Delta}\NormalTok{ v\}\{}\FunctionTok{\textbackslash{}Delta}\NormalTok{ t\} = }\FunctionTok{\textbackslash{}frac}\NormalTok{\{v\_\{}\FunctionTok{\textbackslash{}text}\NormalTok{\{final\}\} {-} v\_\{}\FunctionTok{\textbackslash{}text}\NormalTok{\{start\}\}\}\{}\FunctionTok{\textbackslash{}Delta}\NormalTok{ t\}}
\end{Highlighting}
\end{Shaded}

\begin{itemize}
\tightlist
\item
  \textbf{Positive a} with positive velocity: speeding up.
  \textbf{Negative a} with positive velocity: slowing down. A negative
  acceleration is also called a \emph{deceleration} when the object is
  slowing down.
\item
  \textbf{Zero a} means the velocity is constant. A car at a steady 100
  m/s has a = 0.
\item
  Going around a curve at constant speed is still an acceleration,
  because the \emph{direction} of the velocity changes.
\end{itemize}

\textbf{Getting velocity from a table of positions.} Velocity is change
in position divided by change in time:

\begin{Shaded}
\begin{Highlighting}[]
\NormalTok{v }\FunctionTok{\textbackslash{}approx} \FunctionTok{\textbackslash{}frac}\NormalTok{\{}\FunctionTok{\textbackslash{}Delta}\NormalTok{ x\}\{}\FunctionTok{\textbackslash{}Delta}\NormalTok{ t\}}
\end{Highlighting}
\end{Shaded}

If two position readings are close together in time, this gives the
velocity at the \textbf{midpoint} time between them. To estimate
acceleration, find velocity at two different times, then divide the
change in velocity by the time between those two midpoints.

\textbf{Worked example.} A runner's position is recorded:

{\def\LTcaptype{none} % do not increment counter
\begin{longtable}[]{@{}lllll@{}}
\toprule\noalign{}
Time (s) & 10 & 12 & 50 & 52 \\
\midrule\noalign{}
\endhead
\bottomrule\noalign{}
\endlastfoot
Position (m) & 30 & 42 & 330 & 348 \\
\end{longtable}
}

\begin{enumerate}
\def\labelenumi{\arabic{enumi}.}
\tightlist
\item
  Velocity at t = 11 s: Δx/Δt = (42 − 30) ÷ (12 − 10) = 12 ÷ 2 =
  \textbf{6 m/s}.
\item
  Velocity at t = 51 s: (348 − 330) ÷ (52 − 50) = 18 ÷ 2 = \textbf{9
  m/s}.
\item
  Acceleration from t = 11 s to t = 51 s: (9 − 6) ÷ (51 − 11) = 3 ÷ 40 =
  \textbf{0.075 m/s²}.
\end{enumerate}

\textbf{Worked example (braking).} A car at 30 m/s brakes at −10 m/s².

\begin{itemize}
\tightlist
\item
  Time to stop: t = Δv ÷ a = (0 − 30) ÷ (−10) = \textbf{3 s}.
\item
  Distance: use average velocity. The speed falls steadily from 30 to 0,
  so v\_avg = (30 + 0) ÷ 2 = 15 m/s, and d = 15 × 3 = \textbf{45 m}.
\end{itemize}

\textbf{Try it}

\textbf{B3.1} A cyclist's position is recorded:

{\def\LTcaptype{none} % do not increment counter
\begin{longtable}[]{@{}lllll@{}}
\toprule\noalign{}
Time (s) & 4 & 6 & 14 & 16 \\
\midrule\noalign{}
\endhead
\bottomrule\noalign{}
\endlastfoot
Position (m) & 20 & 32 & 120 & 140 \\
\end{longtable}
}

\begin{enumerate}
\def\labelenumi{(\alph{enumi})}
\tightlist
\item
  Estimate the velocity at t = 5 s. (b) Estimate the velocity at t = 15
  s. (c) Estimate the acceleration from t = 5 s to t = 15 s.
\end{enumerate}

\textbf{B3.2} A car increases its velocity from zero to 60 m/s in 10 s.
What is its acceleration?

\textbf{B3.3} A cart travels at a constant 100 m/s for 10 s. What is its
acceleration during that time?

\textbf{B3.4} A cyclist slows from 12 m/s to 4 m/s in 4 s. Find the
acceleration. What does the sign of your answer tell you?

\textbf{B3.5} An object covers 8 m in the first second, 8 m in the next
second, and 8 m in the third second. What is its acceleration?

\textbf{B3.6} A speeding car at 30 m/s slams on the brakes with a = −10
m/s². (a) How much time passes until it stops? (b) The car's average
velocity while braking is 15 m/s. How far does it travel while stopping?
(c) If instead the car hits a tree and goes from 30 m/s to 0 m/s in 0.05
s, what is the average acceleration in the crash? How many times larger
is it than the braking acceleration?

\textbf{B3.7} A car travels around a circular track at constant speed.
Is it accelerating? Explain.

\textbf{B3.8} ★ A ball rolls to the right at 6 m/s, hits a wall, and
rolls back to the left at 6 m/s, all in 0.2 s. Taking right as positive,
find the average acceleration. (Be careful with the signs of the
velocities.)

\subsection{Part B4: Velocity-time
graphs}\label{part-b4-velocity-time-graphs}

\textbf{Big idea.} A velocity-time graph shows \textbf{how fast and in
which direction} an object moves at each moment. The \textbf{height} of
the line is the velocity, and the \textbf{slope} of the line is the
\textbf{acceleration}:

\begin{Shaded}
\begin{Highlighting}[]
\NormalTok{a = }\FunctionTok{\textbackslash{}text}\NormalTok{\{slope\} = }\FunctionTok{\textbackslash{}frac}\NormalTok{\{}\FunctionTok{\textbackslash{}Delta}\NormalTok{ v\}\{}\FunctionTok{\textbackslash{}Delta}\NormalTok{ t\} = }\FunctionTok{\textbackslash{}frac}\NormalTok{\{}\FunctionTok{\textbackslash{}text}\NormalTok{\{rise\}\}\{}\FunctionTok{\textbackslash{}text}\NormalTok{\{run\}\}}
\end{Highlighting}
\end{Shaded}

\begin{itemize}
\tightlist
\item
  A \textbf{flat} line means \textbf{constant velocity} and zero
  acceleration. The object is at rest only if the flat line is at v = 0.
\item
  A line \textbf{sloping up} means positive acceleration. A line
  \textbf{sloping down} means negative acceleration.
\item
  Where the line \textbf{crosses zero}, the object stops for an instant
  and \textbf{turns around}.
\item
  The object \textbf{speeds up} when the line moves \textbf{away} from
  the zero axis, and \textbf{slows down} when the line moves
  \textbf{toward} the zero axis. In other words, it speeds up when
  velocity and acceleration have the same sign, and slows down when they
  have opposite signs.
\end{itemize}

\textbf{Watch out:} Compare this with Part B2. On a position-time graph
a flat line means \emph{at rest}. On a velocity-time graph a flat line
means \emph{constant velocity}. The same-looking line means something
different, so always check which graph you are reading.

\textbf{Figure 2} shows a cyclist on a straight road. East is the
positive direction.

{[}embedded content: Practice data written for this packet; each dot is
a corner of the graph at whole-number values{]}

\textbf{Worked example.} Find the cyclist's acceleration from t = 0 to t
= 2 s in Figure 2.

\begin{enumerate}
\def\labelenumi{\arabic{enumi}.}
\tightlist
\item
  At t = 0 the velocity is 0. At t = 2 s the velocity is +6 m/s.
\item
  Slope = rise ÷ run = (6 − 0) ÷ (2 − 0) = \textbf{+3 m/s²}.
\item
  Velocity and acceleration are both positive, so the cyclist is
  speeding up toward the east.
\end{enumerate}

\textbf{Try it} (use Figure 2)

\textbf{B4.1} What is the velocity (speed and direction) at t = 4 s?

\textbf{B4.2} What is the velocity at t = 12 s?

\textbf{B4.3} What is the acceleration at t = 4 s?

\textbf{B4.4} What is the acceleration at t = 7 s?

\textbf{B4.5} What is the acceleration at t = 9 s?

\textbf{B4.6} What is the acceleration at t = 15 s?

\textbf{B4.7} At what times is the cyclist momentarily stopped?

\textbf{B4.8} During which time interval is the cyclist moving east?

\textbf{B4.9} During which time intervals is the cyclist (a) speeding
up, (b) slowing down, and (c) moving at constant velocity?

\textbf{B4.10} ★ How far does the cyclist travel during the first 2 s?
(Hint: the velocity rises steadily from 0 to 6 m/s, so use the average
velocity.)

\textbf{B4.11} ★ A student says, ``From 2 s to 6 s the cyclist is
stopped because the line is flat.'' What is wrong with this statement?

\textbf{B4.12} ★ From 8 s to 10 s the acceleration is negative, yet the
cyclist is speeding up. How is that possible?

\subsection{Part B5: Constant acceleration
equations}\label{part-b5-constant-acceleration-equations}

\textbf{Big idea.} When acceleration is \textbf{constant}, velocity
changes by the \textbf{same amount every second}. That makes it possible
to predict the speed and the distance covered at any time. Rearranging a
= Δv ÷ Δt gives the speed equation, and the distance equations come from
using the average velocity.

\begin{Shaded}
\begin{Highlighting}[]
\NormalTok{v = v\_0 + a}\FunctionTok{\textbackslash{},}\NormalTok{t }\FunctionTok{\textbackslash{}qquad\textbackslash{}qquad}\NormalTok{ v\_\{}\FunctionTok{\textbackslash{}text}\NormalTok{\{avg\}\} = }\FunctionTok{\textbackslash{}frac}\NormalTok{\{v\_0 + v\}\{2\} }\FunctionTok{\textbackslash{}qquad\textbackslash{}qquad}\NormalTok{ d = v\_\{}\FunctionTok{\textbackslash{}text}\NormalTok{\{avg\}\}}\FunctionTok{\textbackslash{},}\NormalTok{t}
\end{Highlighting}
\end{Shaded}

If the object \textbf{starts from rest}, then v₀ = 0 and the equations
simplify:

\begin{Shaded}
\begin{Highlighting}[]
\NormalTok{v = a}\FunctionTok{\textbackslash{},}\NormalTok{t }\FunctionTok{\textbackslash{}qquad\textbackslash{}qquad}\NormalTok{ d = }\FunctionTok{\textbackslash{}tfrac}\NormalTok{\{1\}\{2\}}\FunctionTok{\textbackslash{},}\NormalTok{a}\FunctionTok{\textbackslash{},}\NormalTok{t\^{}2}
\end{Highlighting}
\end{Shaded}

\textbf{Which equation do I use?}

{\def\LTcaptype{none} % do not increment counter
\begin{longtable}[]{@{}lll@{}}
\toprule\noalign{}
You know\ldots{} & You want\ldots{} & Use \\
\midrule\noalign{}
\endhead
\bottomrule\noalign{}
\endlastfoot
a and t, starting from rest & Final speed & v = a × t \\
a and t, starting from rest & Distance & d = ½ × a × t² \\
v and a, starting from rest & Time & t = v ÷ a \\
a, t, and a starting speed v₀ & Final speed & v = v₀ + a × t \\
Starting and final speeds, and t & Distance & d = ((v₀ + v) ÷ 2) × t \\
\end{longtable}
}

\textbf{Watch out:} Do not use d = v\_final × t. The object was slower
earlier in the trip, so the distance is smaller than that. Use the
\emph{average} velocity, or use ½at².

\textbf{Worked example.} A car starts from rest and accelerates at 2
m/s².

\begin{enumerate}
\def\labelenumi{\arabic{enumi}.}
\tightlist
\item
  Speed after 3 s: v = at = 2 × 3 = \textbf{6 m/s}.
\item
  Distance in 3 s: d = ½at² = ½ × 2 × 3² = ½ × 2 × 9 = \textbf{9 m}.
\item
  Check with the average: v\_avg = (0 + 6) ÷ 2 = 3 m/s, and d = 3 × 3 =
  9 m. ✓
\item
  Time to reach 30 m/s: t = v ÷ a = 30 ÷ 2 = \textbf{15 s}.
\end{enumerate}

\textbf{Try it}

\textbf{B5.1} A ball starting from rest at the top of a ramp gains 2 m/s
of speed for every second it rolls. What is its acceleration?

\textbf{B5.2} The ramp in B5.1 takes 3 s to roll down. (a) What is the
ball's speed at the bottom? (b) How long is the ramp?

\textbf{B5.3} A car starts from rest and accelerates at 3 m/s² for 5 s.
Find its speed and the distance it travels.

\textbf{B5.4} How long does a car take to go from rest to 24 m/s at 3
m/s²?

\textbf{B5.5} A rocket, starting from rest, accelerates at 50 m/s² for
one minute. What is its speed at the end? (Use 60 s.)

\textbf{B5.6} A cart starts from rest at 5 m/s² and runs for 4 s. Find
the final speed. Then find the distance two ways, with d = ½at² and with
d = v\_avg × t, and show that you get the same answer.

\textbf{B5.7} A car already moving at 20 m/s accelerates at 3 m/s² for 4
s. What is its final speed?

\textbf{B5.8} ★ A student says, ``The car in B5.3 ends at 15 m/s and
drove for 5 s, so it went 15 × 5 = 75 m.'' Find the mistake and give the
correct distance.

\subsection{Part B6: Free fall and
tosses}\label{part-b6-free-fall-and-tosses}

\textbf{Big idea.} An object is in \textbf{free fall} when gravity is
the only force acting on it (we ignore air resistance). Free fall is a
special case of constant acceleration, with a = g ≈ \textbf{10 m/s²
downward}. Five facts cover most questions:

\begin{enumerate}
\def\labelenumi{\arabic{enumi}.}
\tightlist
\item
  The acceleration is \textbf{the same for every object}, heavy or
  light.
\item
  It is the same going \textbf{down, going up, and at the very top}.
  While rising, an object loses 10 m/s of speed every second. While
  falling, it gains 10 m/s every second.
\item
  At the \textbf{top} of a toss, the velocity is momentarily zero but
  the acceleration is still 10 m/s² downward. (If the acceleration were
  zero at the top, the ball would hang there.)
\item
  A ball thrown straight up \textbf{returns with the same speed} it was
  launched with, and the trip up takes \textbf{as long} as the trip
  down.
\item
  Each second a falling object covers a \textbf{greater} distance than
  the second before, since it is speeding up.
\end{enumerate}

\begin{Shaded}
\begin{Highlighting}[]
\NormalTok{v = g}\FunctionTok{\textbackslash{},}\NormalTok{t }\FunctionTok{\textbackslash{}qquad\textbackslash{}qquad}\NormalTok{ d = }\FunctionTok{\textbackslash{}tfrac}\NormalTok{\{1\}\{2\}}\FunctionTok{\textbackslash{},}\NormalTok{g}\FunctionTok{\textbackslash{},}\NormalTok{t\^{}2 = 5}\FunctionTok{\textbackslash{},}\NormalTok{t\^{}2 }\FunctionTok{\textbackslash{}qquad\textbackslash{}qquad} \FunctionTok{\textbackslash{}text}\NormalTok{\{(dropped from rest)\}}
\end{Highlighting}
\end{Shaded}

For an object thrown \textbf{upward} with starting speed v₀, up is
easiest to treat with the symmetry rules:

\begin{Shaded}
\begin{Highlighting}[]
\NormalTok{t\_\{}\FunctionTok{\textbackslash{}text}\NormalTok{\{up\}\} = }\FunctionTok{\textbackslash{}frac}\NormalTok{\{v\_0\}\{g\} }\FunctionTok{\textbackslash{}qquad\textbackslash{}qquad}\NormalTok{ t\_\{}\FunctionTok{\textbackslash{}text}\NormalTok{\{total\}\} = 2}\FunctionTok{\textbackslash{},}\NormalTok{t\_\{}\FunctionTok{\textbackslash{}text}\NormalTok{\{up\}\} }\FunctionTok{\textbackslash{}qquad\textbackslash{}qquad}\NormalTok{ h\_\{}\FunctionTok{\textbackslash{}max}\NormalTok{\} = }\FunctionTok{\textbackslash{}tfrac}\NormalTok{\{1\}\{2\}}\FunctionTok{\textbackslash{},}\NormalTok{g}\FunctionTok{\textbackslash{},}\NormalTok{t\_\{}\FunctionTok{\textbackslash{}text}\NormalTok{\{up\}\}\^{}2}
\end{Highlighting}
\end{Shaded}

(The maximum height is the distance an object would fall from rest in
the time it takes to rise.) You can also use v = v₀ − g × t, where a
positive velocity means upward.

\textbf{Worked example 1: dropped object.} A stone falls from rest down
a mineshaft and hits the bottom after 4 s.

\begin{itemize}
\tightlist
\item
  Speed at the bottom: v = gt = 10 × 4 = \textbf{40 m/s}.
\item
  Depth: d = ½gt² = 5 × 4² = 5 × 16 = \textbf{80 m}.
\end{itemize}

\textbf{Worked example 2: a toss.} A ball is tossed straight up at 30
m/s.

\begin{enumerate}
\def\labelenumi{\arabic{enumi}.}
\tightlist
\item
  Time to rise: t = v₀ ÷ g = 30 ÷ 10 = \textbf{3 s}.
\item
  Total time in the air: 2 × 3 = \textbf{6 s}, and it lands at
  \textbf{30 m/s} downward.
\item
  Maximum height: ½ × 10 × 3² = \textbf{45 m}.
\item
  Velocity after 1 s: v = 30 − 10 × 1 = \textbf{+20 m/s} (up). After 3
  s: \textbf{0} (the top). After 4 s: 30 − 40 = \textbf{−10 m/s} (10 m/s
  downward).
\end{enumerate}

\textbf{Try it}

\textbf{B6.1} A freely falling watermelon falls with constant

\begin{itemize}
\tightlist
\item
  \begin{enumerate}
  \def\labelenumi{\Alph{enumi})}
  \tightlist
  \item
    velocity.
  \end{enumerate}
\item
  \begin{enumerate}
  \def\labelenumi{\Alph{enumi})}
  \setcounter{enumi}{1}
  \tightlist
  \item
    speed.
  \end{enumerate}
\item
  \begin{enumerate}
  \def\labelenumi{\Alph{enumi})}
  \setcounter{enumi}{2}
  \tightlist
  \item
    acceleration.
  \end{enumerate}
\item
  \begin{enumerate}
  \def\labelenumi{\Alph{enumi})}
  \setcounter{enumi}{3}
  \tightlist
  \item
    distance each second.
  \end{enumerate}
\end{itemize}

\textbf{B6.2} A ball is dropped from rest. How fast is it moving after 3
s, and how far has it fallen?

\textbf{B6.3} A stone dropped down a mineshaft takes 6 s to hit the
bottom. How deep is the shaft?

\textbf{B6.4} An apple falls 5 m from a tree. How long does it take to
hit the ground, and how fast is it going when it lands?

\textbf{B6.5} On a distant planet, the acceleration of free fall is 20
m/s². An object is dropped from rest. Find its speed and the distance it
has fallen after 5 s.

\textbf{B6.6} A ball is tossed straight up at 20 m/s. Find (a) its
velocity after 1 s, (b) its velocity after 2 s, (c) its velocity after 3
s, (d) the total time in the air, (e) its acceleration at the very top,
and (f) the maximum height.

\textbf{B6.7} Phil leans over a cliff edge and throws a rock straight up
at 5 m/s. What is the rock's speed 1 s later? And 2 s later? Which
direction is it moving each time?

\textbf{B6.8} How fast must you throw a ball straight up so that it
takes 6 s to return to your hand?

\textbf{B6.9} Phil stands at the edge of a cliff and throws one ball
straight up and another straight down, both at the same speed. Which
ball hits the ground below with the greater speed? Use the idea of
symmetry to explain.

\textbf{B6.10} A basketball player has a hang time of 1 s (total time in
the air). How high does she rise?

\textbf{B6.11} ★ For the ball in B6.6, describe the \emph{velocity-time}
graph. What shape is it, what is its slope, and where does it cross
zero?

\subsection{Part B7: Adding velocity
vectors}\label{part-b7-adding-velocity-vectors}

\textbf{Big idea.} Velocity is a vector, so velocities combine like
forces do. An object's velocity \textbf{relative to the ground} is its
velocity relative to the water or air \textbf{plus} the velocity of the
water or air.

\begin{Shaded}
\begin{Highlighting}[]
\FunctionTok{\textbackslash{}vec}\NormalTok{\{v\}\_\{}\FunctionTok{\textbackslash{}text}\NormalTok{\{ground\}\} = }\FunctionTok{\textbackslash{}vec}\NormalTok{\{v\}\_\{}\FunctionTok{\textbackslash{}text}\NormalTok{\{in the medium\}\} + }\FunctionTok{\textbackslash{}vec}\NormalTok{\{v\}\_\{}\FunctionTok{\textbackslash{}text}\NormalTok{\{medium\}\}}
\end{Highlighting}
\end{Shaded}

There are three cases:

{\def\LTcaptype{none} % do not increment counter
\begin{longtable}[]{@{}
  >{\raggedright\arraybackslash}p{(\linewidth - 4\tabcolsep) * \real{0.3333}}
  >{\raggedright\arraybackslash}p{(\linewidth - 4\tabcolsep) * \real{0.3333}}
  >{\raggedright\arraybackslash}p{(\linewidth - 4\tabcolsep) * \real{0.3333}}@{}}
\toprule\noalign{}
\begin{minipage}[b]{\linewidth}\raggedright
Case
\end{minipage} & \begin{minipage}[b]{\linewidth}\raggedright
Example
\end{minipage} & \begin{minipage}[b]{\linewidth}\raggedright
How to combine
\end{minipage} \\
\midrule\noalign{}
\endhead
\bottomrule\noalign{}
\endlastfoot
Same direction & Plane with a tailwind & Add the speeds \\
Opposite directions & Plane flying into a headwind & Subtract the
speeds; the direction is that of the larger one \\
Right angles & Boat pointed straight across a river current, plane in a
crosswind & R = √(a² + b²) \\
\end{longtable}
}

\textbf{Drawing the vector diagram (tip-to-tail).}

\begin{enumerate}
\def\labelenumi{\arabic{enumi}.}
\tightlist
\item
  Draw the first velocity as an arrow, with its length proportional to
  its speed.
\item
  Start the second arrow at the \textbf{tip} of the first.
\item
  Draw the resultant from the \textbf{tail of the first} arrow to the
  \textbf{tip of the last}.
\item
  Label each arrow with its speed and direction.
\end{enumerate}

For right angles, the three arrows form a right triangle, and the
resultant is its hypotenuse. The resultant always points partly in the
direction of the current or wind. That is why a boat pointed straight
across a river drifts downstream.

\textbf{Worked example.} A boat moves at 6 m/s relative to the water, in
a river that flows east at 2 m/s. Find the boat's velocity relative to
the ground in each case.

\begin{itemize}
\tightlist
\item
  \textbf{Boat points east (with the current):} 6 + 2 = \textbf{8 m/s
  east}.
\item
  \textbf{Boat points west (against the current):} 6 − 2 = \textbf{4 m/s
  west}, since the boat's speed is the larger one.
\item
  \textbf{Boat points north (across the current):} the velocities are at
  right angles, so R = √(6² + 2²) = √40 ≈ \textbf{6.3 m/s}, directed
  mostly north but tilted east.
\end{itemize}

\textbf{Try it}

\textbf{B7.1} An 80 km/h airplane flies in the same direction as a 10
km/h tailwind. What is its speed relative to the ground?

\textbf{B7.2} The same 80 km/h airplane flies into a 10 km/h headwind.
What is its speed relative to the ground?

\textbf{B7.3} An airplane that flies at 100 km/h in still air is caught
in a 100 km/h crosswind. What is its ground speed? (Give your answer to
the nearest whole number.)

\textbf{B7.4} An 80 km/h airplane flies in a 60 km/h crosswind. What is
its ground speed?

\textbf{B7.5} A river 100 m wide flows due south at 1 m/s. A boat that
travels 1 m/s relative to the water is pointed due east and crosses from
the west bank. (a) How long does the crossing take? (b) How far south
does the boat drift during the crossing? (c) What is the boat's speed
relative to the ground, and how far does it actually travel?

\textbf{B7.6} A swimmer swims at 2 m/s relative to the water in a river
that flows east at 1.5 m/s. Find her speed relative to the ground if she
swims (a) east, (b) west, (c) north. Draw a vector diagram for part (c).

\textbf{B7.7} ★ A boat captain wants to land directly across the river
from the starting point, which means the boat must travel straight
across relative to the ground. Should the boat be pointed straight
across, upstream, or downstream? Explain with a vector diagram.

\subsection{Part C: Putting it
together}\label{part-c-putting-it-together}

Chapters 2 and 3 tell one story: \textbf{a net force changes velocity,
and acceleration is how we measure that change.} The table below
collects the connections that show up most often on tests.

{\def\LTcaptype{none} % do not increment counter
\begin{longtable}[]{@{}
  >{\raggedright\arraybackslash}p{(\linewidth - 4\tabcolsep) * \real{0.3333}}
  >{\raggedright\arraybackslash}p{(\linewidth - 4\tabcolsep) * \real{0.3333}}
  >{\raggedright\arraybackslash}p{(\linewidth - 4\tabcolsep) * \real{0.3333}}@{}}
\toprule\noalign{}
\begin{minipage}[b]{\linewidth}\raggedright
If you see\ldots{}
\end{minipage} & \begin{minipage}[b]{\linewidth}\raggedright
Then\ldots{}
\end{minipage} & \begin{minipage}[b]{\linewidth}\raggedright
So the net force is\ldots{}
\end{minipage} \\
\midrule\noalign{}
\endhead
\bottomrule\noalign{}
\endlastfoot
Constant velocity (flat line on a v-t graph, straight line on an x-t
graph) & a = 0 & Zero (equilibrium) \\
Speed changing, direction changing, or both & a ≠ 0 & Not zero \\
Free fall (only gravity acts) & a = 10 m/s² down & Equal to the weight,
pointing down \\
Constant speed around a circle & Direction keeps changing, so a ≠ 0 &
Not zero, pointing toward the center \\
Two perpendicular forces or velocities & Pythagorean theorem for either
& Resultant = √(a² + b²) \\
\end{longtable}
}

\textbf{Worked example.} A ball is tossed straight up. At the very top
of its path, its velocity is zero for an instant. Is the ball in
equilibrium at that moment?

No.~Gravity still pulls it down with a force equal to its weight, and
nothing balances that force. The net force is not zero, so the velocity
is changing (from upward to downward) and the acceleration is 10 m/s²
downward. A velocity of zero for an instant is not the same as having
zero net force.

\textbf{Connecting questions}

\textbf{C1} An insect falls at a constant speed of 3 m/s. (a) What is
its acceleration? (b) What is the net force on it? (c) If its weight is
0.02 N, what is the air resistance on it, and in which direction?

\textbf{C2} A car travels around a circular track at a constant speed.
(a) Is its velocity constant? (b) Is it accelerating? (c) Is the net
force zero? If not, which way does it point?

\textbf{C3} A ball is thrown straight up. At the very top of its path,
what are (a) its velocity, (b) its acceleration, and (c) is it in
equilibrium?

\textbf{C4} A sled's velocity-time graph is a flat line at +4 m/s. What
does that tell you about the net force on the sled? What could the
individual forces be doing?

\textbf{C5} An elevator moves upward at a constant 3 m/s. (a) Compare
the tension in the cable with the elevator's weight. (b) Describe the
shape of its position-time graph and its velocity-time graph.

\textbf{C6} (a) Two ropes pull on a post, one with 30 N toward the north
and one with 40 N toward the east. Find the net force. (b) A boat moves
at 3 m/s north relative to the water in a river that flows 4 m/s east.
Find the boat's speed relative to the ground. (c) What do (a) and (b)
have in common?

\textbf{C7} A hockey puck slides on frictionless ice at 5 m/s. (a) What
is the net force on it? (b) What is its acceleration? (c) Describe its
velocity-time graph and its position-time graph. (d) How far does it
travel in 4 s?

\textbf{C8} ★ A rock is dropped from a high cliff. At first the only
force on it is gravity. (a) What is the net force right after release,
and what is the acceleration? (b) Suppose air resistance grows as the
rock speeds up, and eventually the rock reaches a constant speed. What
must be true about the forces at that point?

\subsection{Mixed practice test}\label{mixed-practice-test}

This is a full practice test. Try it without looking at the earlier
parts, then check your answers. It is worth \textbf{100 points}, and
about 40 minutes is a good target. Use g = 10 m/s² and use Figures 1 and
2 from Parts B2 and B4.

\textbf{Part 1: Multiple choice (4 points each)}

\textbf{T1} A rocket in deep space runs out of fuel. It will

\begin{itemize}
\tightlist
\item
  \begin{enumerate}
  \def\labelenumi{\Alph{enumi})}
  \tightlist
  \item
    keep gaining speed.
  \end{enumerate}
\item
  \begin{enumerate}
  \def\labelenumi{\Alph{enumi})}
  \setcounter{enumi}{1}
  \tightlist
  \item
    slow down gradually until it stops.
  \end{enumerate}
\item
  \begin{enumerate}
  \def\labelenumi{\Alph{enumi})}
  \setcounter{enumi}{2}
  \tightlist
  \item
    continue moving at a constant velocity.
  \end{enumerate}
\item
  \begin{enumerate}
  \def\labelenumi{\Alph{enumi})}
  \setcounter{enumi}{3}
  \tightlist
  \item
    turn around.
  \end{enumerate}
\end{itemize}

\textbf{T2} Which of the following is a vector quantity?

\begin{itemize}
\tightlist
\item
  \begin{enumerate}
  \def\labelenumi{\Alph{enumi})}
  \tightlist
  \item
    Mass
  \end{enumerate}
\item
  \begin{enumerate}
  \def\labelenumi{\Alph{enumi})}
  \setcounter{enumi}{1}
  \tightlist
  \item
    Speed
  \end{enumerate}
\item
  \begin{enumerate}
  \def\labelenumi{\Alph{enumi})}
  \setcounter{enumi}{2}
  \tightlist
  \item
    Time
  \end{enumerate}
\item
  \begin{enumerate}
  \def\labelenumi{\Alph{enumi})}
  \setcounter{enumi}{3}
  \tightlist
  \item
    Force
  \end{enumerate}
\end{itemize}

\textbf{T3} A ball is thrown straight up. At the very top of its path,
its acceleration is

\begin{itemize}
\tightlist
\item
  \begin{enumerate}
  \def\labelenumi{\Alph{enumi})}
  \tightlist
  \item
    zero.
  \end{enumerate}
\item
  \begin{enumerate}
  \def\labelenumi{\Alph{enumi})}
  \setcounter{enumi}{1}
  \tightlist
  \item
    5 m/s² downward.
  \end{enumerate}
\item
  \begin{enumerate}
  \def\labelenumi{\Alph{enumi})}
  \setcounter{enumi}{2}
  \tightlist
  \item
    10 m/s² downward.
  \end{enumerate}
\item
  \begin{enumerate}
  \def\labelenumi{\Alph{enumi})}
  \setcounter{enumi}{3}
  \tightlist
  \item
    10 m/s² upward.
  \end{enumerate}
\end{itemize}

\textbf{T4} A hammock hangs from two ropes tied to two trees. In which
setup is the tension in the ropes greater?

\begin{itemize}
\tightlist
\item
  \begin{enumerate}
  \def\labelenumi{\Alph{enumi})}
  \tightlist
  \item
    Trees close together, so the ropes are steeper
  \end{enumerate}
\item
  \begin{enumerate}
  \def\labelenumi{\Alph{enumi})}
  \setcounter{enumi}{1}
  \tightlist
  \item
    Trees far apart, so the ropes are flatter
  \end{enumerate}
\item
  \begin{enumerate}
  \def\labelenumi{\Alph{enumi})}
  \setcounter{enumi}{2}
  \tightlist
  \item
    The tension is the same in both
  \end{enumerate}
\item
  \begin{enumerate}
  \def\labelenumi{\Alph{enumi})}
  \setcounter{enumi}{3}
  \tightlist
  \item
    There is no tension, because the hammock is at rest
  \end{enumerate}
\end{itemize}

\textbf{T5} A car travels around a circular track at a constant speed.
The net force on the car points

\begin{itemize}
\tightlist
\item
  \begin{enumerate}
  \def\labelenumi{\Alph{enumi})}
  \tightlist
  \item
    toward the center of the circle.
  \end{enumerate}
\item
  \begin{enumerate}
  \def\labelenumi{\Alph{enumi})}
  \setcounter{enumi}{1}
  \tightlist
  \item
    away from the center of the circle.
  \end{enumerate}
\item
  \begin{enumerate}
  \def\labelenumi{\Alph{enumi})}
  \setcounter{enumi}{2}
  \tightlist
  \item
    in the direction the car is moving.
  \end{enumerate}
\item
  \begin{enumerate}
  \def\labelenumi{\Alph{enumi})}
  \setcounter{enumi}{3}
  \tightlist
  \item
    nowhere, because the net force is zero.
  \end{enumerate}
\end{itemize}

\textbf{T6} A freely falling object speeds up each second by about

\begin{itemize}
\tightlist
\item
  \begin{enumerate}
  \def\labelenumi{\Alph{enumi})}
  \tightlist
  \item
    5 m/s.
  \end{enumerate}
\item
  \begin{enumerate}
  \def\labelenumi{\Alph{enumi})}
  \setcounter{enumi}{1}
  \tightlist
  \item
    10 m/s.
  \end{enumerate}
\item
  \begin{enumerate}
  \def\labelenumi{\Alph{enumi})}
  \setcounter{enumi}{2}
  \tightlist
  \item
    20 m/s.
  \end{enumerate}
\item
  \begin{enumerate}
  \def\labelenumi{\Alph{enumi})}
  \setcounter{enumi}{3}
  \tightlist
  \item
    an amount that depends on its mass.
  \end{enumerate}
\end{itemize}

\textbf{T7} An object is in dynamic equilibrium. Which statement must be
true?

\begin{itemize}
\tightlist
\item
  \begin{enumerate}
  \def\labelenumi{\Alph{enumi})}
  \tightlist
  \item
    The object is at rest.
  \end{enumerate}
\item
  \begin{enumerate}
  \def\labelenumi{\Alph{enumi})}
  \setcounter{enumi}{1}
  \tightlist
  \item
    The net force is zero and the velocity is constant.
  \end{enumerate}
\item
  \begin{enumerate}
  \def\labelenumi{\Alph{enumi})}
  \setcounter{enumi}{2}
  \tightlist
  \item
    The net force is zero and the object is speeding up.
  \end{enumerate}
\item
  \begin{enumerate}
  \def\labelenumi{\Alph{enumi})}
  \setcounter{enumi}{3}
  \tightlist
  \item
    No forces act on the object.
  \end{enumerate}
\end{itemize}

\textbf{T8} An airplane flies at 120 km/h through the air in a 50 km/h
crosswind. Its speed relative to the ground is

\begin{itemize}
\tightlist
\item
  \begin{enumerate}
  \def\labelenumi{\Alph{enumi})}
  \tightlist
  \item
    70 km/h.
  \end{enumerate}
\item
  \begin{enumerate}
  \def\labelenumi{\Alph{enumi})}
  \setcounter{enumi}{1}
  \tightlist
  \item
    120 km/h.
  \end{enumerate}
\item
  \begin{enumerate}
  \def\labelenumi{\Alph{enumi})}
  \setcounter{enumi}{2}
  \tightlist
  \item
    130 km/h.
  \end{enumerate}
\item
  \begin{enumerate}
  \def\labelenumi{\Alph{enumi})}
  \setcounter{enumi}{3}
  \tightlist
  \item
    170 km/h.
  \end{enumerate}
\end{itemize}

\textbf{Part 2: Calculations (6 points each, show your work)}

\textbf{T9} An object has these forces on it: 90 N to the right, 30 N to
the left, 20 N up, and 20 N down. (a) Find the net force. (b) Is the
object in equilibrium? (c) What single extra force would put it in
equilibrium?

\textbf{T10} A ball dropped from a bridge hits the water 3 s later. (a)
How fast is it moving when it hits? (b) How high is the bridge?

\textbf{T11} A car starts from rest and accelerates at 4 m/s² for 6 s.
(a) What is its final speed? (b) How far does it travel?

\textbf{T12} A runner's position is recorded:

{\def\LTcaptype{none} % do not increment counter
\begin{longtable}[]{@{}lllll@{}}
\toprule\noalign{}
Time (s) & 9 & 11 & 29 & 31 \\
\midrule\noalign{}
\endhead
\bottomrule\noalign{}
\endlastfoot
Position (m) & 18 & 26 & 146 & 160 \\
\end{longtable}
}

\begin{enumerate}
\def\labelenumi{(\alph{enumi})}
\tightlist
\item
  Estimate the velocity at t = 10 s. (b) Estimate the velocity at t = 30
  s. (c) Estimate the acceleration from t = 10 s to t = 30 s.
\end{enumerate}

\textbf{Part 3: Graphs (8 points each)}

\textbf{T13} Use Figure 1. (a) What is the student's position at t = 14
s? (b) What is the student's velocity at t = 10 s? (c) During which
interval(s) is the student at rest?

\textbf{T14} Use Figure 2. (a) What is the cyclist's acceleration at t =
7 s? (b) What is the cyclist's velocity at t = 12 s? (c) During which
interval(s) is the cyclist speeding up?

\textbf{Part 4: Explanations (8 points each)}

\textbf{T15} Using Newton's first law, explain why a passenger lurches
forward when the car they are riding in brakes suddenly.

\textbf{T16} A student says, ``At the top of its flight, a ball tossed
straight up is in equilibrium, because it is momentarily at rest.'' Is
she right? Explain.

\textbf{Part 5: Vectors (12 points)}

\textbf{T17} A swimmer swims east at 1.5 m/s relative to the water in a
river that flows north at 2.0 m/s. (a) Draw a tip-to-tail vector
diagram. (4 points) (b) Find the swimmer's speed relative to the ground.
(4 points) (c) What would the swimmer's velocity relative to the ground
be if the swimmer instead swam south at 1.5 m/s relative to the water?
(4 points)

\subsection{Answer key}\label{answer-key}

Check your work against these answers after you try each question. If
your answer differs, find the step where your thinking went a different
way.

\textbf{Part A1: Inertia}

\begin{itemize}
\tightlist
\item
  \textbf{A1.1} A.
\item
  \textbf{A1.2} C. No force is needed to keep an object moving at
  constant velocity.
\item
  \textbf{A1.3} B. The rock moves in a straight line in the direction it
  was moving when the string broke.
\item
  \textbf{A1.4} The bowling ball. It has more mass, so it has more
  inertia and is harder to start, stop, or turn.
\item
  \textbf{A1.5} Friction (and a little air resistance) pushes backward
  on the puck, so there is a net force that slows it. On frictionless
  ice there would be no net force, and the puck would slide at constant
  velocity forever.
\item
  \textbf{A1.6} The coin shares the train's forward velocity, and
  nothing pushes it forward or backward, so it keeps that velocity and
  lands in your hand. If the train speeds up, the coin keeps its old
  forward velocity while the train moves ahead, so it lands behind you.
\end{itemize}

\textbf{Part A2: Net force}

\begin{itemize}
\tightlist
\item
  \textbf{A2.1} 25 − 10 = 15 N toward the south.
\item
  \textbf{A2.2} 30 + 20 − 10 = 40 N to the right.
\item
  \textbf{A2.3} Up-down: 100 − 60 = 40 N up. Left-right: 35 − 35 = 0.
  Net force is 40 N up.
\item
  \textbf{A2.4} False. Zero net force means the velocity does not
  change, so the object could be at rest or moving at constant velocity.
\item
  \textbf{A2.5} The effect of a force depends on both how strong it is
  and which way it points. Two forces with the same magnitude but
  different directions have different effects.
\item
  \textbf{A2.6} (a) Both forces point in the same direction. (b) The
  forces point in opposite directions.
\item
  \textbf{A2.7} The cart slows down. If the net force kept acting, it
  would stop and then speed up toward the left.
\end{itemize}

\textbf{Part A3: Equilibrium}

\begin{itemize}
\tightlist
\item
  \textbf{A3.1} A (static equilibrium).
\item
  \textbf{A3.2} Yes, dynamic equilibrium, because its velocity is
  constant. The net force is 0 N.
\item
  \textbf{A3.3} The known forces add to +120 − 45 = 75 N to the right,
  so the unknown force is \textbf{75 N to the left}.
\item
  \textbf{A3.4} The net force is 0 N. Air resistance is \textbf{600 N
  upward}, equal to her weight.
\item
  \textbf{A3.5} \textbf{2000 N backward}, equal and opposite to the
  engine force. It is dynamic equilibrium.
\item
  \textbf{A3.6} Not necessarily. If the box slides at constant velocity,
  the net force is zero (the push balances friction). A net force is
  needed only if the box's velocity is changing.
\item
  \textbf{A3.7} The net force is zero in both directions, so the object
  is in equilibrium. But equilibrium includes constant-velocity motion,
  so it could be at rest or moving at constant velocity.
\end{itemize}

\textbf{Part A4: Support force, tension, force diagrams}

\begin{itemize}
\tightlist
\item
  \textbf{A4.1} 10 N, upward.
\item
  \textbf{A4.2} Weight = 30 × 10 = 300 N. Support force = 300 N, upward.
\item
  \textbf{A4.3} 250 N on each scale.
\item
  \textbf{A4.4} 1300 + 400 = 1700 N.
\item
  \textbf{A4.5} Net force on the rope is 0 N. Tension is 300 N.
\item
  \textbf{A4.6} (a) Friction is 240 N to the left, to balance the 240 N
  rightward part of the tension. (b) Up: 650 + 150 = 800 N, which equals
  the 800 N weight down.
\item
  \textbf{A4.7} No.~The person is still at rest, so the support force
  still equals the weight.
\item
  \textbf{A4.8} The forward part of the tension is now greater than the
  friction force, so there is a net force forward and the sled is not in
  equilibrium.
\end{itemize}

\textbf{Part A5: Angles, tension, right-angle forces}

\begin{itemize}
\tightlist
\item
  \textbf{A5.1} B. The flatter the rope, the greater the tension.
\item
  \textbf{A5.2} 400 N each.
\item
  \textbf{A5.3} Greater than 400 N. Only part of each tension points
  upward, so each rope has to pull harder for the upward parts to add up
  to 800 N.
\item
  \textbf{A5.4} The ropes become flatter, so the angle gets smaller and
  the tension gets larger.
\item
  \textbf{A5.5} √(6² + 8²) = √100 = 10 N. (The 6-8-10 triangle.)
\item
  \textbf{A5.6} About 14 N.
\item
  \textbf{A5.7} √(5² + 12²) = √169 = 13 N.
\item
  \textbf{A5.8} A perfectly flat rope would have tension pointing only
  sideways, with no upward part to support the shirt's weight. The
  tension would have to be infinite, so the rope always sags.
\end{itemize}

\textbf{Part B1: Speed, velocity, relative motion}

\begin{itemize}
\tightlist
\item
  \textbf{B1.1} Instantaneous speed.
\item
  \textbf{B1.2} 400 ÷ 80 = 5 m/s.
\item
  \textbf{B1.3} 2 min = 120 s, so d = 6 × 120 = 720 m.
\item
  \textbf{B1.4} 45 min = 0.75 h, so 60 ÷ 0.75 = 80 km/h.
\item
  \textbf{B1.5} Velocity and acceleration.
\item
  \textbf{B1.6} 11 m/s if she runs in the same direction as the train,
  or 5 m/s if she runs in the opposite direction.
\item
  \textbf{B1.7} 0 m/s.
\item
  \textbf{B1.8} The speed is constant, but the velocity is not, because
  the direction keeps changing.
\end{itemize}

\textbf{Part B2: Position-time graph (Figure 1)}

\begin{itemize}
\tightlist
\item
  \textbf{B2.1} 8 m.
\item
  \textbf{B2.2} 0 m (the doorway).
\item
  \textbf{B2.3} 0 m/s (flat line, so at rest).
\item
  \textbf{B2.4} Slope = (−8 − 8) ÷ (12 − 8) = −16 ÷ 4 = \textbf{−4 m/s}:
  4 m/s in the negative direction (toward the west end).
\item
  \textbf{B2.5} Slope from 16 s to 20 s = (−4 − (−8)) ÷ 4 = \textbf{+1
  m/s}.
\item
  \textbf{B2.6} From 4 s to 8 s and from 12 s to 16 s.
\item
  \textbf{B2.7} From 0 s to 10 s (after the start).
\item
  \textbf{B2.8} From 0 s to 4 s and from 16 s to 20 s.
\item
  \textbf{B2.9} From 8 s to 12 s, at 4 m/s (the steepest line).
\item
  \textbf{B2.10} No.~The position is zero at that instant, but the slope
  is −4 m/s, so the student is moving.
\item
  \textbf{B2.11} Distance = 8 + 16 + 4 = \textbf{28 m} (the total length
  of the walk). The final position is \textbf{−4 m}, which is 4 m toward
  the west end from the doorway. Distance adds up every step, while
  position only says where the student ended up.
\end{itemize}

\textbf{Part B3: Tables and acceleration}

\begin{itemize}
\tightlist
\item
  \textbf{B3.1} (a) (32 − 20) ÷ 2 = 6 m/s. (b) (140 − 120) ÷ 2 = 10 m/s.
  (c) (10 − 6) ÷ (15 − 5) = 0.4 m/s².
\item
  \textbf{B3.2} 60 ÷ 10 = 6 m/s².
\item
  \textbf{B3.3} 0 m/s², because the velocity is constant.
\item
  \textbf{B3.4} (4 − 12) ÷ 4 = \textbf{−2 m/s²}. The negative sign means
  velocity is decreasing: the cyclist is slowing down while moving in
  the positive direction.
\item
  \textbf{B3.5} 0 m/s². It moves 8 m every second, which is a constant
  velocity of 8 m/s.
\item
  \textbf{B3.6} (a) t = (0 − 30) ÷ (−10) = \textbf{3 s}. (b) d = 15 × 3
  = \textbf{45 m}. (c) a = (0 − 30) ÷ 0.05 = \textbf{−600 m/s²}, which
  is 60 times the braking acceleration.
\item
  \textbf{B3.7} Yes. Its velocity changes because the direction changes,
  even though the speed does not.
\item
  \textbf{B3.8} a = (−6 − 6) ÷ 0.2 = \textbf{−60 m/s²}.
\end{itemize}

\textbf{Part B4: Velocity-time graph (Figure 2)}

\begin{itemize}
\tightlist
\item
  \textbf{B4.1} +6 m/s (6 m/s east).
\item
  \textbf{B4.2} −4 m/s (4 m/s west).
\item
  \textbf{B4.3} 0 m/s² (flat line).
\item
  \textbf{B4.4} Slope from 6 s to 8 s = (0 − 6) ÷ 2 = \textbf{−3 m/s²}.
\item
  \textbf{B4.5} (−4 − 0) ÷ 2 = \textbf{−2 m/s²}.
\item
  \textbf{B4.6} (0 − (−4)) ÷ 2 = \textbf{+2 m/s²}.
\item
  \textbf{B4.7} At t = 0, t = 8 s, and t = 16 s (where the line is at v
  = 0).
\item
  \textbf{B4.8} From 0 s to 8 s (velocity is positive).
\item
  \textbf{B4.9} (a) Speeding up: 0 to 2 s and 8 to 10 s. (b) Slowing
  down: 6 to 8 s and 14 to 16 s. (c) Constant velocity: 2 to 6 s and 10
  to 14 s.
\item
  \textbf{B4.10} Average velocity = (0 + 6) ÷ 2 = 3 m/s, so d = 3 × 2 =
  \textbf{6 m}.
\item
  \textbf{B4.11} A flat line on a velocity-time graph means constant
  velocity (+6 m/s east), not rest. The cyclist is stopped only where
  the line is at zero.
\item
  \textbf{B4.12} The velocity is negative (moving west) and the
  acceleration is negative too. When they have the same sign, the speed
  increases. On the graph, the line moves farther from zero.
\end{itemize}

\textbf{Part B5: Constant acceleration equations}

\begin{itemize}
\tightlist
\item
  \textbf{B5.1} 2 m/s².
\item
  \textbf{B5.2} (a) v = 2 × 3 = 6 m/s. (b) d = ½ × 2 × 3² = 9 m.
\item
  \textbf{B5.3} v = 3 × 5 = \textbf{15 m/s}. d = ½ × 3 × 5² =
  \textbf{37.5 m}.
\item
  \textbf{B5.4} t = v ÷ a = 24 ÷ 3 = 8 \textbf{s}.
\item
  \textbf{B5.5} v = 50 × 60 = \textbf{3000 m/s}.
\item
  \textbf{B5.6} v = 5 × 4 = 20 m/s. d = ½ × 5 × 4² = 40 m. With
  averages: v\_avg = (0 + 20) ÷ 2 = 10 m/s, and d = 10 × 4 = 40 m. Both
  agree.
\item
  \textbf{B5.7} v = 20 + 3 × 4 = \textbf{32 m/s}.
\item
  \textbf{B5.8} The student used the final speed for the whole trip, but
  the car was slower earlier. The average speed is (0 + 15) ÷ 2 = 7.5
  m/s, so d = 7.5 × 5 = \textbf{37.5 m}.
\end{itemize}

\textbf{Part B6: Free fall}

\begin{itemize}
\tightlist
\item
  \textbf{B6.1} C (acceleration).
\item
  \textbf{B6.2} v = 10 × 3 = 30 m/s. d = 5 × 3² = 45 m.
\item
  \textbf{B6.3} d = 5 × 6² = 5 × 36 = 180 m.
\item
  \textbf{B6.4} 5 = 5t², so t = 1 s. v = 10 × 1 = 10 m/s.
\item
  \textbf{B6.5} v = 20 × 5 = 100 m/s. d = ½ × 20 × 5² = 250 m.
\item
  \textbf{B6.6} (a) +10 m/s (10 m/s up). (b) 0 m/s (the top). (c) −10
  m/s (10 m/s down). (d) 4 s. (e) 10 m/s² downward. (f) ½ × 10 × 2² = 20
  m.
\item
  \textbf{B6.7} After 1 s: v = 5 − 10 = −5, so \textbf{5 m/s downward}
  (the rock rose for half a second, then started falling). After 2 s: v
  = 5 − 20 = −15, so \textbf{15 m/s downward}.
\item
  \textbf{B6.8} It must rise for 3 s, so v₀ = 10 × 3 = \textbf{30 m/s}.
\item
  \textbf{B6.9} Both hit with the same speed. The ball thrown up rises,
  stops, and comes back down past the cliff edge with the same speed it
  was thrown with, going downward. From there it is exactly like the
  ball thrown down.
\item
  \textbf{B6.10} She rises for 0.5 s, so h = 5 × 0.5² = \textbf{1.25 m}.
\item
  \textbf{B6.11} It is a straight line sloping downward with slope −10
  m/s². It starts at +20 m/s at t = 0, crosses zero at t = 2 s (the top
  of the path), and ends at −20 m/s at t = 4 s.
\end{itemize}

\textbf{Part B7: Adding velocity vectors}

\begin{itemize}
\tightlist
\item
  \textbf{B7.1} 90 km/h.
\item
  \textbf{B7.2} 70 km/h.
\item
  \textbf{B7.3} √(100² + 100²) ≈ \textbf{141 km/h}.
\item
  \textbf{B7.4} √(80² + 60²) = \textbf{100 km/h}.
\item
  \textbf{B7.5} (a) 100 m ÷ 1 m/s = 100 s. (b) 1 m/s × 100 s = 100 m
  south. (c) Speed = √(1² + 1²) ≈ 1.4 m/s, and the distance traveled is
  about 1.4 × 100 ≈ 141 m.
\item
  \textbf{B7.6} (a) 2 + 1.5 = \textbf{3.5 m/s east}. (b) 2 − 1.5 =
  \textbf{0.5 m/s west}. (c) √(2² + 1.5²) = √6.25 = \textbf{2.5 m/s},
  toward the north and tilted east.
\item
  \textbf{B7.7} Point the boat \textbf{upstream}. The upstream part of
  the boat's velocity cancels the current, so the resultant points
  straight across. In the tip-to-tail diagram, the boat's arrow is
  angled upstream, the current's arrow points downstream from its tip,
  and the resultant goes straight across.
\end{itemize}

\textbf{Part C: Connecting questions}

\begin{itemize}
\tightlist
\item
  \textbf{C1} (a) 0 m/s². (b) 0 N. (c) 0.02 N upward.
\item
  \textbf{C2} (a) No, the direction changes. (b) Yes. (c) No.~The net
  force points toward the center of the circle.
\item
  \textbf{C3} (a) 0 m/s. (b) 10 m/s² downward. (c) No, because the
  weight is not balanced.
\item
  \textbf{C4} The net force is zero (dynamic equilibrium). The forward
  pull on the sled equals the backward friction force.
\item
  \textbf{C5} (a) The tension equals the weight. (b) The position-time
  graph is a straight line sloping upward (slope 3 m/s). The
  velocity-time graph is a flat line at +3 m/s.
\item
  \textbf{C6} (a) √(30² + 40²) = 50 N. (b) √(3² + 4²) = 5 m/s. (c) Both
  combine two perpendicular vectors, so both use the Pythagorean
  theorem.
\item
  \textbf{C7} (a) 0 N. (b) 0 m/s². (c) The velocity-time graph is a flat
  line at 5 m/s. The position-time graph is a straight line sloping
  upward with slope 5 m/s. (d) 5 × 4 = 20 m.
\item
  \textbf{C8} (a) The net force equals the rock's weight, pointing down,
  so the acceleration is 10 m/s² downward. (b) Air resistance has grown
  until it equals the weight. The net force is zero, the acceleration is
  zero, and the rock is in dynamic equilibrium.
\end{itemize}

\textbf{Mixed practice test}

\begin{itemize}
\tightlist
\item
  \textbf{T1} C. \textbf{T2} D. \textbf{T3} C. \textbf{T4} B.
  \textbf{T5} A. \textbf{T6} B. \textbf{T7} B. \textbf{T8} C (√(50² +
  120²) = 130).
\item
  \textbf{T9} (a) 60 N to the right (the up and down forces cancel). (b)
  No, the net force is not zero. (c) 60 N to the left.
\item
  \textbf{T10} (a) v = 10 × 3 = 30 m/s. (b) d = 5 × 3² = 45 m.
\item
  \textbf{T11} (a) v = 4 × 6 = 24 m/s. (b) d = ½ × 4 × 6² = 72 m.
\item
  \textbf{T12} (a) (26 − 18) ÷ 2 = 4 m/s. (b) (160 − 146) ÷ 2 = 7 m/s.
  (c) (7 − 4) ÷ (30 − 10) = 0.15 m/s².
\item
  \textbf{T13} (a) −8 m. (b) −4 m/s (4 m/s in the negative direction).
  (c) 4 to 8 s and 12 to 16 s.
\item
  \textbf{T14} (a) −3 m/s². (b) −4 m/s (4 m/s west). (c) 0 to 2 s and 8
  to 10 s.
\item
  \textbf{T15} The passenger is moving forward with the car. When the
  brakes slow the car, nothing slows the passenger at the same moment,
  so by inertia the passenger keeps moving forward at the original
  velocity. Relative to the slowing car, the passenger seems to lurch
  forward.
\item
  \textbf{T16} She is not right. At the top, the velocity is zero for an
  instant, but gravity still acts on the ball and nothing balances it.
  The net force is the weight, the acceleration is 10 m/s² downward, and
  the velocity is changing. Equilibrium requires zero net force.
\item
  \textbf{T17} (a) Draw an arrow east for 1.5 m/s, then start an arrow
  north for 2.0 m/s at the tip of the first. The resultant goes from the
  start of the first arrow to the tip of the second. (b) √(1.5² + 2.0²)
  = √6.25 = \textbf{2.5 m/s}. (c) 2.0 north − 1.5 south = \textbf{0.5
  m/s north}.
\end{itemize}

\end{document}
